\NormalFont\ShortTitle{Jó}
{\MT Jó

\par }\ChapOne{1}{\PP \VerseOne{1}Havia um homem na terra de Uz, cujo nome era Jó; e este homem era íntegro e correto, temente a Deus, e que se afastava do mal.
\VS{2}E nasceram-lhe sete filhos e três filhas.
\VS{3}E seu patrimônio era sete mil ovelhas, três mil camelos, quinhentas juntas de bois, e quinhentas jumentas; ele também tinha muitíssimos servos, de maneira que este homem era o maior de todos do oriente.
\VS{4}E seus filhos iam nas casas uns dos outros para fazerem banquetes, cada um em seu dia; e mandavam convidar as suas três irmãs, para que comessem e bebessem com eles.
\VS{5}E acontecia que, acabando-se o revezamento dos dias de banquetes, Jó enviava e os santificava, e se levantava de madrugada para apresentar holocaustos{\ADD{conforme}} o número de todos eles. Pois Jó dizia: Talvez meus filhos tenham pecado, e tenham amaldiçoado a Deus em seus corações. Assim Jó fazia todos aqueles dias.
\VS{6}E em certo dia, os filhos de Deus
\FTNT{*}{{\FR 1:6 }
filhos de Deus
i.e., os anjos
} vieram para se apresentarem diante do SENHOR, e Satanás também veio entre eles.
\VS{7}Então o SENHOR disse a Satanás: De onde vens? E Satanás respondeu ao SENHOR, dizendo: De rodear a terra, e de passear por ela.
\VS{8}E o SENHOR disse a Satanás: Tendes visto meu servo Jó? Pois ninguém há na terra semelhante a ele, homem íntegro e correto, temente a Deus, e que se afasta de mal.
\VS{9}Então Satanás respondeu ao SENHOR, dizendo: Por acaso Jó teme a Deus em troca de nada?
\VS{10}Por acaso tu não puseste uma cerca ao redor dele, de sua casa, e de tudo quanto ele tem? Tu abençoaste o trabalho de suas mãos, e seu patrimônio
\FTNT{†}{{\FR 1:10 }
patrimônio
trad. alt. gado
} tem crescido sobre a terra.
\VS{11}Mas estende agora tua mão, e toca em tudo quanto ele tem; e
{\ADD{verás}} se ele não te amaldiçoa em tua face.
\VS{12}E o SENHOR disse a Satanás: Eis que tudo quanto ele tem está em tua mão; somente não estendas tua mão contra ele. E Satanás saiu de diante do SENHOR.
\VS{13}E sucedeu um dia que seus filhos e filhas estavam comendo e bebendo vinho na casa de seu irmão primogênito,
\VS{14}Que veio um mensageiro a Jó, que disse: Enquanto os bois estavam arando, e as jumentas se alimentando perto deles,
\VS{15}Eis que os sabeus atacaram, e os tomaram, e feriram os servos a fio de espada; somente eu escapei para te trazer a notícia.
\VS{16}Enquanto este ainda estava falando, veio outro que disse: Fogo de Deus caiu do céu, que incendiou as ovelhas entre os servos, e os consumiu; somente eu escapei para te trazer-te a notícia.
\VS{17}Enquanto este ainda estava falando, veio outro que disse: Os caldeus formaram três tropas, e atacaram os camelos, e os tomaram, e feriram os servos a fio de espada; somente eu escapei para te trazer a notícia.
\VS{18}Enquanto este ainda estava falando, veio outro que disse: Teus filhos e tuas filhas estavam comendo e bebendo vinho em casa de seu irmão primogênito,
\VS{19}E eis que veio um grande vento do deserto, e atingiu os quatro cantos da casa, que caiu sobre os jovens, e morreram; somente eu escapei para te trazer a notícia.
\VS{20}Então Jó se levantou, rasgou sua capa, rapou sua cabeça, e caindo na terra, adorou,
\VS{21}E disse: Nu saí do ventre de minha mãe, e nu para lá voltarei. O SENHOR deu, e o SENHOR tomou; bendito seja o nome do SENHOR.
\VS{22}Em tudo isto Jó não pecou, nem atribuiu a Deus falta alguma.

\par }\Chap{2}{\PP \VerseOne{1}E veio outro dia em que os filhos de Deus vieram para se apresentarem diante do SENHOR, e e Satanás também veio entre eles para se apresentar diante do SENHOR.
\VS{2}Então o SENHOR disse a Satanás: De onde vens? E Satanás respondeu ao SENHOR, e dizendo: De rodear a terra, e de passear por ela.
\VS{3}E o SENHOR disse a Satanás: Tendes visto meu servo Jó? Pois ninguém há semelhante a ele na terra, homem íntegro e correto, temente a Deus e que se afasta do mal; e que ainda mantém sua integridade, mesmo depois de teres me incitado contra ele, para o arruinar sem causa.
\VS{4}Então Satanás respondeu ao SENHOR, dizendo: Pele por pele, e tudo quanto o homem tem, dará por sua vida.
\VS{5}Mas estende agora tua mão, e toca em seus ossos e sua carne, e
{\ADD{verás}} se ele não te amaldiçoa em tua face.
\VS{6}E o SENHOR disse a Satanás: Eis que ele está em tua mão; mas preserva sua vida.
\VS{7}Então Satanás saiu de diante do SENHOR, e feriu a Jó de chagas malignas desde a planta de seus pés até a topo de sua cabeça.
\VS{8}E
{\ADD{Jó}} tomou um caco para se raspar com ele, e ficou sentado no meio da cinza.
\VS{9}Então sua mulher lhe disse: Ainda manténs a tua integridade? Amaldiçoa a Deus, e morre.
\VS{10}Porém ele lhe disse: Tu falas como uma tola. Receberíamos o bem de Deus, e o mal não receberíamos? Em tudo isto Jó não pecou com seus lábios.
\VS{11}Quando três amigos de Jó, Elifaz temanita, Bildade suíta, e Zofar naamita, ouviram todo este mal que lhe tinha vindo sobre ele, vieram cada um de seu lugar; porque haviam combinado de juntamente virem para se condoerem dele, e o consolarem.
\VS{12}Eles, quando levantaram seus olhos de longe, não o reconheceram; e choraram em alta voz; e cada um deles rasgou sua capa, e espalharam pó ao ar sobre suas cabeças.
\VS{13}Assim se sentaram com ele na terra durante sete dias e sete noites, e nenhum lhe falava palavra alguma, pois viam que a dor era muito grande.

\par }\Chap{3}{\PP \VerseOne{1}Depois disto Jó abriu sua boca, e amaldiçoou seu dia.
\VS{2}Pois Jó respondeu, e disse:
\VS{3}Pereça o dia em que nasci, e a noite
{\ADD{em que}} se disse: Um homem foi concebido.
\VS{4}Torne-se aquele dia em trevas; Deus não lhe dê atenção desde acima, nem claridade brilhe sobre ele.
\VS{5}Reivindiquem-no para si trevas e sombra de morte; nuvens habitem sobre ele; a escuridão do dia o espante.
\VS{6}Tome a escuridão aquela noite; não seja contada entre os dias do ano, nem faça parte do número dos meses.
\VS{7}Ah se aquela noite fosse solitária, e música de alegria não viesse a ela!
\VS{8}Amaldiçoem-na os que amaldiçoam o dia, os que se preparam para levantar seu pranto.
\VS{9}Escureçam-se as estrelas de sua manhã; espere a luz, e não venha, e as pálpebras não vejam o amanhecer;
\VS{10}Pois não fechou as portas do ventre onde eu estava, nem escondeu de meus olhos o sofrimento.
\VS{11}Por que eu não morri desde a madre, ou perdi a vida ao sair do ventre?
\VS{12}Por que joelhos me receberam? E por que seios me amamentaram?
\VS{13}Pois agora eu jazeria e repousaria; dormiria, e então haveria repouso para mim;
\VS{14}Com os reis e os conselheiros da terra, que edificavam para si os desertos;
\VS{15}Ou com os príncipes que tinham ouro, que enchiam suas casas de prata.
\VS{16}Ou
{\ADD{por que}} não fui como um aborto oculto, como as crianças que nunca viram a luz?
\VS{17}Ali os maus deixam de perturbar, e ali repousam os cansados de forças.
\VS{18}Ali os prisioneiros juntamente repousam;
{\ADD{e}} não ouvem a voz do opressor.
\VS{19}Ali estão o pequeno e o grande; e o servo livre
{\ADD{está}} de seu senhor.
\VS{20}Por que se dá luz ao sofredor, e vida aos amargos de alma,
\VS{21}Que esperam a morte, e ela não chega, e que a buscam mais que tesouros;
\VS{22}Que saltam de alegram e ficam contentes quando acham a sepultura?
\VS{23}{\ADD{E também}} ao homem cujo caminho é oculto, e a quem Deus
{\ADD{o}} encobriu?
\VS{24}Pois antes do meu pão vem meu suspiro; e meus gemidos correm como águas.
\VS{25}Pois aquilo eu temia tanto veio a mim, e aquilo que tinha medo me aconteceu.
\VS{26}Não tenho tido descanso, nem tranquilidade, nem repouso; mas perturbação veio sobre mim.

\par }\Chap{4}{\PP \VerseOne{1}Então Elifaz o temanita respondeu, dizendo:
\VS{2}Se tentarmos falar contigo, ficarás incomodado? Mas quem poderia deter as palavras?
\VS{3}Eis que tu ensinavas a muitos, e fortalecias as mãos fracas;
\VS{4}Tuas palavras levantavam aos que tropeçavam, e fortificavas os joelhos que desfaleciam.
\VS{5}Mas agora
{\ADD{isso}} que aconteceu contigo, tu te cansas; e quando
{\ADD{isso}} te tocou, te perturbas.
\VS{6}Por acaso não era o teu temor
{\ADD{a Deus}} a tua confiança, e a integridade dos teus caminhos tua esperança?
\VS{7}Lembra-te agora, qual foi o inocente que pereceu? E onde os corretos foram destruídos?
\VS{8}Como eu tenho visto, os que lavram injustiça e semeiam opressão colhem o mesmo.
\VS{9}Com o sopro de Deus eles perecem, e pelo vento de sua ira
\FTNT{*}{{\FR 4:9 }
sua ira
trad. alt. suas narinas
} são consumidos.
\VS{10}O rugido do leão, a voz do leão feroz, e os dentes dos leões jovens são quebrantados.
\VS{11}O leão velho perece por falta de presa, e os filhotes da leoa se dispersam.
\VS{12}Uma palavra me foi dita em segredo, e meu ouvidos escutaram um sussurro dela.
\VS{13}Em imaginações de visões noturnas, quando o sono profundo cai sobre os homens,
\VS{14}Espanto e tremor vieram sobre mim, que espantou todos os meus ossos.
\VS{15}Então um vento
\FTNT{†}{{\FR 4:15 }
vento
trad. alt. espírito
} passou por diante de mim, que fez arrepiar os pelos de minha carne.
\VS{16}Ele parou, mas não reconheci sua feição; uma figura estava diante de meus olhos, e ouvi uma voz quieta,
{\ADD{que dizia}} :
\VS{17}Seria o ser humano mais justo que Deus? Seria o homem mais puro que seu Criador?
\VS{18}Visto que ele não confia em seus servos, e considera seus anjos como loucos,
\VS{19}Quanto mais naqueles que habitam em casas de lodo, cujo fundamento está no pó, e são esmagáveis como a traça!
\VS{20}Desde a manhã até a tarde são despedaçados, e perecem sempre, sem que ninguém perceba.
\VS{21}Por acaso sua excelência não se perde com eles mesmos? Eles morrem sem sabedoria.

\par }\Chap{5}{\PP \VerseOne{1}Clama agora! Haverá alguém que te responda? E a qual dos santos te voltarás?
\VS{2}Pois a ira acaba com o louco, e o zelo mata o tolo.
\VS{3}Eu vi ao louco lançar raízes, porém logo amaldiçoei sua habitação.
\VS{4}Seus filhos estarão longe da salvação; na porta são despedaçados, e não há quem os livre.
\VS{5}O faminto devora sua colheita, e a tira até dentre os espinhos; e o assaltante
\FTNT{*}{{\FR 5:5 }
assaltante – obscuro
trad. alt. sedento
} traga sua riqueza.
\VS{6}Porque a aflição não procede do pó da terra, nem a opressão brota do chão.
\VS{7}Mas o ser humano nasce para a opressão, assim como as faíscas das brasas se levantam a voar.
\VS{8}Porém eu buscaria a Deus, e a ele confiaria minha causa;
\VS{9}{\ADD{Pois}} ele é o que faz coisas grandiosas e incompreensíveis, e inúmeras maravilhas.
\VS{10}Ele é o que dá a chuva sobre a face da terra, e envia águas sobre os campos.
\VS{11}Ele põe os humildes em lugares altos, para que os sofredores sejam postos em segurança.
\VS{12}Ele frustra os planos dos astutos, para que suas mãos nada consigam executar.
\VS{13}Ele prende aos sábios em sua própria astúcia; para que o conselho dos perversos seja derrubado.
\VS{14}De dia eles se encontram com as trevas, e ao meio-dia andam apalpando como de noite.
\VS{15}Porém livra ao necessitado da espada de suas bocas, e da mão do violento.
\VS{16}Pois ele é esperança para o necessitado, e a injustiça tapa sua boca.
\VS{17}Eis que bem-aventurado é o homem a quem Deus corrige; portanto não rejeites o castigo do Todo-Poderoso.
\VS{18}Pois ele faz a chaga, mas também põe o curativo; ele fere, mas suas mãos curam.
\VS{19}Em seis angústias ele te livrará, e em sete o mal não te tocará.
\VS{20}Na fome ele te livrará da morte, e na guerra
{\ADD{livrará}} do poder da espada.
\VS{21}Do açoite da língua estarás encoberto; e não temerás a destruição quando ela vier.
\VS{22}Tu rirás da destruição e da fome, e não temerás os animais da terra.
\VS{23}Pois até com as pedras do campo terás teu pacto, e os animais do campo serão pacíficos contigo.
\VS{24}E saberás que há paz em tua tenda; e visitarás tua habitação, e não falharás.
\FTNT{†}{{\FR 5:24 }
falharás
trad. alt. acharás coisa alguma em falta
}
\VS{25}Também saberás que tua semente se multiplicará, e teus descendentes serão como a erva da terra.
\VS{26}Na velhice virás à sepultura, como o amontoado de trigo que se recolhe a seu tempo.
\VS{27}Eis que é isto o que temos constatado, e assim é; ouve-o, e pensa nisso tu para teu
{\ADD{bem}} .

\par }\Chap{6}{\PP \VerseOne{1}Mas Jó respondeu, dizendo:
\VS{2}Oh se pesassem justamente minha aflição, e meu tormento juntamente fosse posto em uma balança!
\VS{3}Pois na verdade seria mais pesada que a areia dos mares; por isso minhas palavras têm sido impulsivas.
\VS{4}Porque as flechas do Todo-Poderoso estão em mim, cujo veneno meu espírito bebe; e temores de Deus me atacam.
\VS{5}Por acaso o asno selvagem zurra junto à erva, ou o boi berra junto a seu pasto?
\VS{6}Por acaso se come o insípido sem sal? Ou há gosto na clara do ovo?
\VS{7}Minha alma se recusa tocar
{\ADD{essas coisas}} ,que são para mim como comida detestável.
\VS{8}Ah se meu pedido fosse realizado, e se Deus
{\ADD{me}} desse o que espero!
\VS{9}Que Deus me destruísse; ele soltasse sua mão, e acabasse comigo!
\VS{10}Isto ainda seria meu consolo, um alívio em meio ao tormento que não
{\ADD{me}} poupa; pois eu não tenho escondido as palavras do Santo.
\VS{11}Qual é minha força para que eu espere? E qual meu fim, para que eu prolongue minha vida?
\VS{12}É, por acaso, a minha força a força de pedras? Minha carne é de bronze?
\VS{13}Tenho eu como ajudar a mim mesmo, se todo auxílio
\FTNT{*}{{\FR 6:13 }
auxílio
obscuro
} me foi tirado?
\VS{14}Ao aflito, seus amigos deviam ser misericordiosos, mesmo se ele tivesse abandonado o temor ao Todo-Poderoso.
\VS{15}Meus irmãos foram traiçoeiros comigo, como ribeiro, como correntes de águas que transbordam,
\VS{16}Que estão escurecidas pelo gelo, e nelas se esconde a neve;
\VS{17}Que no tempo do calor se secam e, ao se aquecerem, desaparecem de seu lugar;
\VS{18}Os cursos de seus caminhos se desviam; vão se minguando, e perecem.
\VS{19}As caravanas de Temã as veem; os viajantes de Sabá esperam por elas.
\VS{20}Foram envergonhados por aquilo em que confiavam; e ao chegarem ali, ficaram desapontados.
\VS{21}Agora, vós vos tornastes semelhantes a elas; pois vistes o terror, e temestes.
\VS{22}Por acaso eu disse: Trazei-me
{\ADD{algo}} ? Ou: Dai presente a mim de vossa riqueza?
\VS{23}Ou: Livrai-me da mão do opressor? Ou: Resgatai-me das mãos dos violentos?
\VS{24}Ensinai-me, e eu
{\ADD{me}} calarei; e fazei-me entender em que errei.
\VS{25}Como são fortes as palavras de boa razão! Mas o que vossa repreensão reprova?
\VS{26}Pretendeis repreender palavras, sendo que os argumentos do desesperado são como o vento?
\VS{27}De fato vós lançaríeis
{\ADD{sortes}} sobre o órfão, e venderíeis vosso amigo.
\VS{28}Agora, pois, disponde-vos a olhar para mim; e
{\ADD{vede}} se eu minto diante de vós.
\VS{29}Mudai de opinião, pois, e não haja perversidade; mudai de opinião, pois minha justiça continua.
\VS{30}Há perversidade em minha língua? Não poderia meu paladar discernir as coisas más?

\par }\Chap{7}{\PP \VerseOne{1}Por acaso o ser humano não tem um trabalho duro sobre a terra, e não são seus dias como os dias de um assalariado?
\VS{2}Como o servo suspira pela sombra, e como o assalariado espera por seu pagamento,
\VS{3}Assim também me deram por herança meses inúteis, e me prepararam noites de sofrimento.
\VS{4}Quando eu me deito, pergunto: Quando me levantarei? Mas a noite se prolonga, e me canso de me virar
{\ADD{na cama}} até o amanhecer.
\VS{5}Minha carne está coberta de vermes e de crostas de pó; meu pele está rachada e horrível.
\VS{6}Meus dias são mais rápidos que a lançadeira do tecelão, e perecem sem esperança.
\VS{7}Lembra-te que minha vida é um sopro; meus olhos não voltarão a ver o bem.
\VS{8}Os olhos dos que me veem não me verão mais; teus olhos estarão sobre mim, porém deixarei de existir.
\VS{9}A nuvem se esvaece, e passa; assim também quem desce ao Xeol
\FTNT{*}{{\FR 7:9 }
Xeol é o lugar dos mortos
} nunca voltará a subir.
\VS{10}Nunca mais voltará à sua casa, nem seu lugar o conhecerá.
\VS{11}Por isso eu não calarei minha boca; falarei na angústia do meu espírito, e me queixarei na amargura de minha alma.
\VS{12}Por acaso sou eu o mar, ou um monstro marinho, para que me ponhas guarda?
\VS{13}Quando eu digo: Minha cama me consolará; meu leito aliviará minhas queixa,
\VS{14}Então tu me espantas com sonhos, e me assombras com visões.
\VS{15}Por isso minha alma preferia a asfixia
{\ADD{e}} a morte, mais que meus ossos.
\VS{16}Odeio
{\ADD{a minha vida}} ; não viverei para sempre; deixa-me, pois que meus dias são inúteis.
\VS{17}O que é o ser humano, para que tanto o estimes, e ponhas sobre ele teu coração,
\VS{18}E o visites a cada manhã, e a cada momento o proves?
\VS{19}Até quando não me deixarás, nem me liberarás até que eu engula minha saliva?
\VS{20}Se pequei, o eu que te fiz, ó Guarda dos homens? Por que me fizeste de alvo de dardos, para que eu seja pesado para mim mesmo?
\VS{21}E por que não perdoas minha transgressão, e tiras minha maldade? Porque agora dormirei no pó, e me buscarás de manhã, porém não mais existirei.

\par }\Chap{8}{\PP \VerseOne{1}Então Bildade, o suíta, respondeu, dizendo:
\VS{2}Até quando falarás tais coisas, e as palavras de tua boca serão como um vento impetuoso?
\VS{3}Por acaso Deus perverteria o direito, ou o Todo-Poderoso perverteria a justiça?
\VS{4}Se teus filhos pecaram contra ele, ele também os entregou ao castigo por sua transgressão.
\FTNT{*}{{\FR 8:4 }
castigo por seu pecado
lit. poder de sua transgressão
}
\VS{5}Se tu buscares a Deus com empenho, e pedires misericórdia ao Todo-Poderoso;
\VS{6}Se fores puro e correto, certamente logo ele se levantará em teu favor, e restaurará a morada de tua justiça.
\VS{7}Ainda que teu princípio seja pequeno, o teu fim será muito grandioso.
\VS{8}Pois pergunta agora à geração passada, e considera o que seus pais descobriram.
\VS{9}Pois nós somos de ontem e nada sabemos, pois nossos dias sobre a terra são como a sombra.
\VS{10}Por acaso eles não te ensinarão, e te dirão, e falarão palavras de seu coração?
\VS{11}Pode o papiro crescer sem lodo? Ou pode o junco ficar maior sem água?
\VS{12}Estando ele ainda verde, sem ter sido cortado, ainda assim se seca antes de toda erva.
\VS{13}Assim são os caminhos de todos os que esquecem de Deus; e a esperança do corrupto perecerá;
\VS{14}Sua esperança será frustrada, e sua confiança será como a teia de aranha.
\VS{15}Ele se apoiará em sua casa, mas ela não ficará firme; ele se apegará a ela, mas ela não ficará de pé.
\VS{16}Ele está bem regado diante do sol, e seus ramos brotam por cima de sua horta;
\VS{17}Suas raízes se entrelaçam junto à fonte, olhando para o pedregal.
\VS{18}Se lhe arrancarem de seu lugar, este o negará,
{\ADD{dizendo}} : Nunca te vi.
\VS{19}Eis que este é o prazer de seu caminho; e do solo outros brotarão.
\VS{20}Eis que Deus não rejeita ao íntegro, nem segura pela mão aos malfeitores.
\VS{21}Ainda ele encherá tua boca de riso, e teus lábios de júbilo.
\VS{22}Os que te odeiam se vestirão de vergonha, e nunca mais haverá tenda de perversos.

\par }\Chap{9}{\PP \VerseOne{1}Mas Jó respondeu, dizendo:
\VS{2}Na verdade sei que é assim; mas como pode o ser humano ser justo diante de Deus?
\VS{3}Ainda se quisesse disputar com ele, não conseguiria lhe responder uma coisa sequer em mil.
\VS{4}Ele é sábio de coração, e poderoso em forças. Quem se endureceu contra ele, e teve paz?
\VS{5}Ele transporta as montanhas sem que o saibam; e as transtorna em seu furor.
\VS{6}Ele remove a terra de seu lugar, e faz suas colunas tremerem.
\VS{7}Ele dá ordem ao sol, e ele não brilha; e sela as estrelas.
\VS{8}Ele é o que sozinho estende os céus, e anda sobre as alturas do mar.
\VS{9}Ele é o que fez a Usra, o Órion, as Plêiades, e as constelações
\FTNT{*}{{\FR 9:9 }
constelações
lit. recâmaras
} do sul.
\VS{10}Ele é o que faz coisas grandes e incompreensíveis, e inúmeras maravilhas.
\VS{11}Eis que ele passa diante de mim, sem que eu não o veja; ele passará diante de mim, sem que eu saiba.
\VS{12}Eis que, quando ele toma, quem pode lhe impedir? Quem poderá lhe dizer: O que estás fazendo?
\VS{13}Deus não reverterá sua ira, e debaixo dele se encurvam os assistentes de Raabe.
\FTNT{†}{{\FR 9:13 }
assistentes de Raabe
obscuro – possivelmente monstros marinhos – tradicionalmente traduzido como assistentes soberbos
}
\VS{14}Como poderia eu lhe responder, e escolher minhas palavras contra ele?
\VS{15}A ele, ainda que eu fosse justo, não lhe responderia; a meu juiz pediria misericórdia.
\VS{16}Ainda que eu lhe chamasse, e ele respondesse, mesmo assim não creria que ele tivesse dado ouvidos à minha voz.
\VS{17}Pois ele tem me quebrantado com tempestade, e multiplicado minhas feridas sem causa.
\VS{18}Ele não me permite respirar; em vez disso, me farta de amarguras.
\VS{19}Quanto às forças, eis que ele é forte; e quanto ao juízo,
{\ADD{ele diria}} : Quem me convocará?
\VS{20}Ainda que eu seja justo, minha boca me condenaria; se eu fosse inocente, então ela me declararia perverso.
\VS{21}Mesmo se eu for inocente, não estimo minha alma; desprezo minha vida.
\VS{22}É tudo a mesma coisa; por isso digo: ele consome ao inocente e ao perverso.
\VS{23}Quando o açoite mata de repente, ele ri do desespero dos inocentes.
\VS{24}A terra está entregue nas mãos dos perversos. Ele cobre o rosto de seus juízes. Se não é ele, então quem é?
\VS{25}Meus dias foram mais rápidos que um homem que corre; fugiram, e não viram o bem.
\VS{26}Passaram como barcos de papiro, como a águia que se lança à comida.
\VS{27}Se disser: Esquecerei minha queixa, mudarei o aspecto do meu rosto, e sorrirei,
\VS{28}{\ADD{Ainda}} teria pavor de todas as minhas dores;
{\ADD{pois}} sei que não me terás por inocente.
\VS{29}Se eu já estou condenado, então para que eu sofreria em vão?
\VS{30}Ainda que me lave com água de neve, e limpe minhas mãos com sabão,
\VS{31}Então me submergirias no fosso, e minhas próprias vestes me abominariam.
\VS{32}Pois ele não é homem como eu, para que eu lhe responda, e venhamos juntamente a juízo.
\VS{33}Não há entre nós árbitro que ponha sua mão sobre nós ambos,
\VS{34}Tire de sobre mim sua vara, e seu terror não me espante.
\VS{35}{\ADD{Então}} eu falaria, e não teria medo dele. Pois não está sendo assim comigo.

\par }\Chap{10}{\PP \VerseOne{1}Minha alma está cansada de minha vida. Darei liberdade à minha queixa sobre mim; falarei com amargura de minha alma.
\VS{2}Direi a Deus: Não me condenes; faz-me saber por que brigas comigo.
\VS{3}{\ADD{Parece}} -te bem que
{\ADD{me}} oprimas, que rejeites o trabalho de tuas mãos, e favoreças
\FTNT{*}{{\FR 10:3 }
favoreças
brilhes, com o sentido de sorrias
} o conselho dos perversos?
\VS{4}Tens tu olhos de carne? Vês tu como o ser humano vê?
\VS{5}São teus dias como os dias do ser humano, ou teus anos como os anos do homem,
\VS{6}Para que investigues minha perversidade, e pesquises meu pecado?
\VS{7}Tu sabes que eu não sou mau; todavia ninguém há que
{\ADD{me}} livre de tua mão.
\VS{8}Tuas mãos me fizeram e me formaram por completo; porém agora tu me destróis.
\VS{9}Por favor, lembra-te que me preparaste como o barro; e me farás voltar ao pó da terra.
\VS{10}Por acaso não me derramaste como o leite, e como o queijo me coalhaste?
\VS{11}De pele e carne tu me vestiste; e de ossos e nervos tu me teceste.
\VS{12}Vida e misericórdia me concedeste, e teu cuidado guardou meu espírito.
\VS{13}Porém estas coisas escondeste em teu coração; eu sei que isto esteve contigo:
\VS{14}Se eu pecar, tu me observarás, e não absolverás minha culpa.
\VS{15}Se eu for perverso, ai de mim! Mesmo se eu for justo, não levantarei minha cabeça; estou farto de desonra, e de ver minha aflição.
\VS{16}Se
{\ADD{minha cabeça}} se exaltar, tu me caças como um leão feroz, e voltas a fazer em coisas extraordinárias contra mim.
\VS{17}Renovas tuas testemunhas contra mim, e multiplicas tua ira sobre mim; combates vêm sucessivamente contra mim.
\VS{18}Por que me tiraste da madre?
{\ADD{Bom seria}} se eu não tivesse respirado, e nenhum olho me visse!
\VS{19}Teria sido como se nunca tivesse existido, e desde o ventre
{\ADD{materno}} seria levado à sepultura.
\VS{20}Por acaso não são poucos os meus dias? Cessa
{\ADD{pois}} e deixa-me, para que eu tenha um pouco de alívio,
\VS{21}Antes que eu me vá para não voltar, à terra da escuridão e da sombra de morte;
\VS{22}Terra escura ao extremo, tenebrosa, sombra de morte, sem ordem alguma, onde a luz é como a escuridão.

\par }\Chap{11}{\PP \VerseOne{1}Então Zofar, o naamita, respondeu, dizendo:
\VS{2}Por acaso a multidão de palavras não seria respondida? E o homem falador teria razão?
\VS{3}Por acaso tuas palavras tolas faria as pessoas se calarem? E zombarias tu, e ninguém te envergonharia?
\VS{4}Pois disseste: Minha doutrina é pura, e eu sou limpo diante de teus olhos.
\VS{5}Mas na verdade, queria eu que Deus falasse, e abrisse seus lábios contra ti,
\VS{6}E te fizesse saber os segredos da sabedoria, porque o verdadeiro conhecimento tem dois lados; por isso sabe tu que Deus tem te castigado menos que
{\ADD{mereces}} por tua perversidade.
\VS{7}Podes tu compreender os mistérios de Deus? Chegarás tu à perfeição do Todo-Poderoso?
\VS{8}{\ADD{Sua sabedoria}} é mais alta que os céus; que tu poderás fazer? E mais profunda que o Xeol;
\FTNT{*}{{\FR 11:8 }
Xeol é o lugar dos mortos
} que tu poderás conhecer?
\VS{9}Sua medida é mais comprida que a terra, e mais larga que o mar.
\VS{10}Se ele passar, e prender, ou se ajuntar
{\ADD{para o julgamento}} ,quem poderá o impedir?
\VS{11}Pois ele conhece as pessoas vãs, e vê a maldade; por acaso ele não a consideraria?
\VS{12}O homem tolo se tornará entendido quando do asno selvagem nascer um humano.
\VS{13}Se tu preparares o teu coração, e estenderes a ele tuas mãos;
\VS{14}Se alguma maldade houver em tua mão, lança-a de ti, e não deixes a injustiça habitar em tuas tendas;
\VS{15}Então levantarás teu rosto sem mácula; estarás firme, e não temerás;
\VS{16}E esquecerás teu sofrimento,
{\ADD{ou}} lembrarás dele como de águas que já passaram;
\VS{17}E
{\ADD{tua}} vida será mais clara que o meio-dia; ainda que haja trevas, tu serás como o amanhecer.
\VS{18}E serás confiante, porque haverá esperança; olharás em redor,
{\ADD{e}} repousarás seguro.
\VS{19}E te deitarás, e ninguém te causará medo; e muitos suplicarão a ti.
\VS{20}Porém os olhos dos maus se enfraquecerão, e o refúgio deles perecerá; a esperança deles será a morte.
\FTNT{†}{{\FR 11:20 }
a morte
lit. o expirar da alma
}

\par }\Chap{12}{\PP \VerseOne{1}Porém Jó respondeu, dizendo:
\VS{2}Verdadeiramente vós sois o povo; e convosco morrerá a sabedoria.
\VS{3}Também eu tenho entendimento
\FTNT{*}{{\FR 12:3 }
entendimento
lit. coração – também v. 24
} como vós; e não sou inferior a vós; e quem há que não saiba coisas como essas?
\VS{4}Eu sou o motivo de riso de meus amigos, eu que invocava a Deus, e ele me respondia; o justo e íntegro serve de riso.
\VS{5}Na opinião de quem está descansado, a desgraça é desprezada,
{\ADD{como se}} estivesse preparada aos que cujos pés escorregam.
\VS{6}As tendas dos ladrões têm descanso, e os que irritam a Deus estão seguros; os que trazem
{\ADD{seu}} deus em suas mãos.
\VS{7}Verdadeiramente pergunta agora aos animais, que eles te ensinarão; e às aves dos céus, que elas te explicarão;
\VS{8}Ou fala com a terra, que ela te ensinará; até os peixes do mar te contarão.
\VS{9}Quem entre todas estas coisas não entende que a mão do SENHOR faz isto?
\VS{10}Em sua mão está a alma de tudo quanto vive, e o espírito
\FTNT{†}{{\FR 12:10 }
espírito
i.e. fôlego
} de toda carne humana.
\VS{11}Por acaso o ouvido não distingue as palavras, e o paladar prova as comidas?
\VS{12}Nos velhos está o conhecimento, e na longa idade
\FTNT{‡}{{\FR 12:12 }
longa idade
lit. longura de dias
} o entendimento.
\VS{13}Com
{\ADD{Deus}}
\FTNT{§}{{\FR 12:13 }
[Deus]
lit. ele
} está a sabedoria e a força; o conselho e o entendimento lhe pertencem.
\VS{14}Eis que o que ele derruba não pode ser reconstruído; e ninguém pode libertar o homem a quem ele aprisiona.
\VS{15}Eis que,
{\ADD{quando}} ele detém as águas, elas se secam;
{\ADD{quando}} ele as deixa sair, elas transtornam a terra.
\VS{16}Com ele está a força e a sabedoria; Seu é o que erra, e o que faz errar.
\VS{17}Ele leva os conselheiros despojados, e faz os juízes enlouquecerem.
\VS{18}Ele solta a atadura dos reis, e ata um cinto a seus lombos.
\VS{19}Ele leva os sacerdotes despojados, e transtorna os poderosos.
\VS{20}Ele tira a fala daqueles a quem os outros confiam, e tira o juízo dos anciãos.
\VS{21}Ele derrama menosprezo sobre os príncipes, e afrouxa o cinto dos fortes.
\VS{22}Ele revela as profundezas das trevas, e traz a sombra de morte à luz.
\VS{23}Ele multiplica as nações, e ele as destrói; ele dispersa as nações, e as reúne.
\VS{24}Ele tira o entendimento dos líderes
\FTNT{*}{{\FR 12:24 }
líderes
lit. cabeças
} do povo da terra, e os faz vaguear pelos desertos, sem caminho.
\VS{25}Nas trevas andam apalpando, sem terem luz; e os faz cambalear como a bêbados.

\par }\Chap{13}{\PP \VerseOne{1}Eis que meus olhos têm visto tudo
{\ADD{isto}} ; meus ouvidos o ouviram, e entenderam.
\VS{2}Assim como vós o sabeis, eu também o sei; não sou inferior a vós.
\VS{3}Mas eu falarei com o Todo-Poderoso, e quero me defender para com Deus.
\VS{4}Pois na verdade vós sois inventores de mentiras; todos vós sois médicos inúteis.
\VS{5}Bom seria se vos calásseis por completo, pois seria sabedoria de vossa parte.
\VS{6}Ouvi agora meu argumento, e prestai atenção aos argumentos de meus lábios.
\VS{7}Por acaso falareis perversidade por Deus, e por ele falareis engano?
\VS{8}Fareis acepção de sua pessoa? Brigareis em defesa de Deus?
\VS{9}Seria bom
{\ADD{para vós}} se ele vos investigasse? Enganareis a ele como se engana a algum homem?
\VS{10}Certamente ele vos repreenderá, se em oculto fizerdes acepção de pessoas.
\VS{11}Por acaso a majestade dele não vos espantará? E o temor dele não cairá sobre sobre vós?
\VS{12}Vossos conceitos são provérbios de cinzas; vossas defesas são como defesas de lama.
\VS{13}Calai-vos diante de mim, e eu falarei; e venha sobre mim o que vier.
\VS{14}Por que tiraria eu minha carne com meus dentes, e poria minha alma em minha mão?
\VS{15}Eis que, ainda que ele me mate, nele esperarei; porém defenderei meus caminhos diante dele.
\VS{16}Ele mesmo será minha salvação; pois o hipócrita não virá perante ele.
\VS{17}Ouvi com atenção minhas palavras, e com vossos ouvidos minha declaração.
\VS{18}Eis que já tenho preparado minha causa; sei que serei considerado justo.
\VS{19}Quem é o que brigará comigo? Pois então eu me calaria e morreria.
\FTNT{*}{{\FR 13:19 }
morreria
lit. expiraria
}
\VS{20}Somente duas coisas não faças comigo; então eu não me esconderei de teu rosto:
\VS{21}Afasta tua mão de sobre mim, e teu terror não me espante.
\VS{22}Chama, e eu responderei; ou eu falarei, e tu me responde.
\VS{23}Quantas culpas e pecados eu tenho? Faze-me saber minha transgressão e meu pecado.
\VS{24}Por que escondes teu rosto, e me consideras teu inimigo?
\VS{25}Por acaso quebrarás a folha arrebatada
{\ADD{pelo vento}} ? E perseguirás a palha seca?
\VS{26}Por que escreves contra mim amarguras, e me fazes herdar as transgressões de minha juventude?
\VS{27}Também pões meus pés no tronco, e observas todos os meus caminhos. Tu pões limites às solas dos meus pés.
\VS{28}Eu
\FTNT{†}{{\FR 13:28 }
Eu
lit. ele
} me consumo como a podridão, como uma roupa que a traça rói.

\par }\Chap{14}{\PP \VerseOne{1}O homem, nascido de mulher, é curto de dias, e farto de inquietação;
\VS{2}Ele sai como uma flor, e é cortado; foge como a sombra, e não permanece.
\VS{3}Contudo sobre este abres teus olhos, e me trazes a juízo contigo.
\VS{4}Quem tirará algo puro do imundo? Ninguém.
\VS{5}Visto que seus dias já estão determinados, e contigo está o número de seus meses, tu lhe puseste limites,
{\ADD{dos quais}} ele não passará.
\VS{6}Desvia-te dele, para que ele tenha repouso; até que, como o empregado, complete seu dia.
\VS{7}Porque há
{\ADD{ainda}} alguma esperança para a árvore que, se cortada, ainda se renove, e seus renovos não cessem.
\VS{8}Ainda que sua raiz se envelheça na terra, e seu tronco morra no solo,
\VS{9}Ao cheiro das águas ela brotará, e dará ramos como uma planta nova.
\VS{10}Porém o homem morre, e se abate; depois de expirar, onde ele está?
\VS{11}As águas se vão do lago, e o rio se esgota, e se seca.
\VS{12}Assim o homem se deita, e não se levanta; até que não haja mais céus, eles não despertarão, nem se erguerão de seu sono.
\VS{13}Ah, se tu me escondesses no Xeol,
\FTNT{*}{{\FR 14:13 }
Xeol é o lugar dos mortos
} e me ocultasses até que a tua ira passasse, se me pusesses um limite de tempo, e te lembrasses de mim!
\VS{14}Se o homem morrer, voltará a viver? Todos os dias de meu combate esperarei, até que venha minha dispensa.
\VS{15}Tu
{\ADD{me}} chamarás, e eu te responderei; e te afeiçoarás à obra de tuas mãos.
\VS{16}Pois então tu contarias meus passos, e não ficarias vigiando meu pecado.
\VS{17}Minha transgressão estaria selada numa bolsa, e tu encobririas minhas perversidades.
\VS{18}E assim como a montanha cai e é destruída, e a rocha muda de seu lugar,
\VS{19}E a água desgasta as pedras, e as enxurradas levam o pó da terra, assim também tu fazes perecer a esperança do homem.
\VS{20}Sempre prevaleces contra ele, e ele passa; tu mudas o aspecto de seu rosto, e o despedes.
\VS{21}Se seus filhos vierem a ter honra, ele não saberá; se forem humilhados, ele não perceberá.
\VS{22}Ele apenas sente as dores em sua própria carne, e lamenta por sua própria alma.

\par }\Chap{15}{\PP \VerseOne{1}Então Elifaz, o temanita, respondeu, dizendo:
\VS{2}Por acaso o sábio dará como resposta vão conhecimento, e encherá seu ventre de vento oriental?
\VS{3}Repreenderá com palavras que nada servem, e com argumentos que de nada aproveitam?
\VS{4}Porém tu destróis o temor, e menosprezas a oração diante de Deus.
\VS{5}Pois tua perversidade conduz tua boca, e tu escolheste a língua dos astutos.
\VS{6}Tua boca te condena, e não eu; e teus lábios dão testemunho contra ti.
\VS{7}Por acaso foste tu o primeiro ser humano a nascer? Ou foste gerado antes dos morros?
\VS{8}Ouviste tu o segredo de Deus? Reténs tu
{\ADD{apenas}} contigo a sabedoria?
\VS{9}O que tu sabes que nós não saibamos?
{\ADD{O que}} tu entendes que não tenhamos
{\ADD{entendido}} ?
\VS{10}Entre nós também há os que tenham cabelos grisalhos, também há os que são muito mais idosos que teu pai.
\VS{11}Por acaso as consolações de Deus te são poucas? As mansas palavras voltadas a ti?
\VS{12}Por que o teu coração te arrebata, e por que centelham teus olhos,
\VS{13}Para que vires teu espírito contra Deus, e deixes sair
{\ADD{tais}} tais palavras de tua boca?
\VS{14}O que é o homem, para que seja puro? E o nascido de mulher, para que seja justo?
\VS{15}Eis que
{\ADD{Deus}} não confia em seus santos, nem os céus são puros diante de seus olhos;
\VS{16}Quanto menos o homem, abominável e corrupto, que bebe a maldade como água?
\VS{17}Escuta-me; eu te mostrarei; eu te contarei o que vi.
\VS{18}(O que os sábios contaram, o que não foi encoberto por seus pais,
\VS{19}A somente os quais a terra foi dada, e estranho nenhum passou por meio deles):
\VS{20}Todos os dias do perverso são sofrimento para si, o número de anos reservados ao opressor.
\VS{21}Ruídos de horrores estão em seus ouvidos; até na paz lhe sobrevém o assolador.
\VS{22}Ele não crê que voltará da escuridão; ao contrário, a espada o espera.
\VS{23}Anda vagueando por comida, onde quer que ela esteja. Ele sabe que o dia das trevas está prestes a acontecer.
\FTNT{*}{{\FR 15:23 }
prestes a acontecer
lit. preparado em sua mão
}
\VS{24}Angústia e aflição o assombram,
{\ADD{e}} prevalecem contra ele como um rei preparado para a batalha;
\VS{25}Porque ele estendeu sua mão contra Deus, e se embraveceu contra o Todo-Poderoso,
\VS{26}Corre contra ele com
{\ADD{dureza}} de pescoço, e como seus escudos grossos e levantados.
\VS{27}Porque cobriu seu rosto com sua gordura, e engordou as laterais de seu corpo.
\VS{28}E habitou em cidades desoladas cidades, em casas desabitadas; que estavam prestes a desmoronar.
\VS{29}Ele não enriquecerá, nem seu patrimônio subsistirá, nem suas riquezas se estenderão pela terra.
\VS{30}Não escapará das trevas; a chama secará seus ramos, e ao sopro de sua boca desparecerá.
\VS{31}Não confie ele na ilusão para ser enganado; pois a sua recompensa será nada.
\VS{32}Não sendo ainda seu tempo, ela se cumprirá; e seu ramo não florescerá.
\VS{33}Sacudirá suas uvas antes de amadurecerem como a vide, e derramará sua flor como a oliveira.
\VS{34}Pois a ajuntamento dos hipócritas será estéril, e fogo consumirá as tendas do suborno.
\VS{35}Eles concebem a maldade, e dão à luz a perversidade; e o ventre deles prepara enganos.

\par }\Chap{16}{\PP \VerseOne{1}Porém Jó respondeu, dizendo:
\VS{2}Ouvi muitas coisas como estas; todos vós sois consoladores miseráveis.
\VS{3}Por acaso terão fim as palavras de vento? Ou o que é que te provoca a responderes?
\VS{4}Também eu poderia falar como vós, se vossa alma estivesse no lugar da minha alma; eu poderia amontoar palavras contra vós, e contra vós sacudir minha cabeça.
\VS{5}Porém eu vos confortaria com minha boca, e a consolação de meus lábios serviria para aliviar.
\VS{6}Ainda que eu fale, minha dor não cessa; e se eu me calar, em que me alivio?
\VS{7}Na verdade agora ele me tornou exausto; tu assolaste toda a minha companhia.
\VS{8}Testemunha
{\ADD{disto é}} que já me enrugaste; e minha magreza já se levanta contra mim para em meu rosto dar testemunho
{\ADD{contra mim}} .
\VS{9}Sua ira me despedaça, e ele me odeia; range seus dentes contra mim; meu adversário aguça seus olhos contra mim.
\VS{10}Abrem sua boca contra mim; com desprezo esbofeteiam meu rosto, e todos se ajuntam contra mim.
\VS{11}Deus me entregou ao perverso, e me fez cair nas mãos dos malignos.
\VS{12}Tranquilo eu estava, porém ele me quebrantou; e pegou-me pelo pescoço, e me despedaçou; e fez de mim seu alvo de pontaria.
\VS{13}Seus flecheiros me cercaram-me, partiu meus rins, e não
{\ADD{me}} poupou; meu fel derramou em terra.
\VS{14}Quebrantou-me de quebrantamento sobre quebrantamento; correu contra mim como um guerreiro.
\VS{15}Costurei saco sobre minha pele, e revolvi minha cabeça no pó.
\VS{16}Meu rosto está vermelho de choro, e minhas pálpebras estão escurecidas ao extremo;
\VS{17}Apesar de não haver injustiça em minhas mãos, e de minha oração ser pura.
\VS{18}Ó terra! Não cubras o meu sangue, e não haja lugar para meu clamor!
\VS{19}Eis que mesmo agora minha testemunha está nos céus, e meu defensor nas alturas.
\VS{20}Meus amigos zombam de mim,
{\ADD{mas}} meus olhos estão derramando para Deus.
\VS{21}Ah, se
{\ADD{fosse possível}} defender a causa com Deus em favor do homem, como o filho do homem em favor de seu amigo!
\VS{22}Pois poucos anos restam, e seguirei o caminho
{\ADD{por onde}} não voltarei.

\par }\Chap{17}{\PP \VerseOne{1}Meu espírito está arruinado, meus dias vão se extinguindo, e a sepultura já etá pronta para mim.
\VS{2}Comigo há ninguém além de zombadores, e meus olhos são obrigados a ficar diante de suas provocações.
\VS{3}Concede-me, por favor, uma garantia para comigo; quem
{\ADD{outro}} há que me dê a mão?
\VS{4}Pois aos corações deles tu encobriste do entendimento; portanto não os exaltarás.
\VS{5}Aquele que denuncia a seus amigos em proveito próprio, também os olhos de seus filhos desfalecerão.
\VS{6}Ele tem me posto por ditado de povos, e em meu rosto é onde eles cospem.
\VS{7}Por isso meus olhos se escureceram de mágoa, e todos os membros de meu corpo são como a sombra.
\VS{8}Os íntegros pasmarão sobre isto, e o inocente se levantará contra o hipócrita.
\VS{9}E o justo prosseguirá seu caminho, e o puro de mãos crescerá em força.
\VS{10}Mas, na verdade, voltai-vos todos vós, e vinde agora, pois sábio nenhum acharei entre vós.
\VS{11}Meus dias se passaram, meus pensamentos foram arrancados, os desejos do meu coração.
\VS{12}Tornaram a noite em dia; a luz se encurta por causa das trevas.
\VS{13}Se eu esperar, o Xeol
\FTNT{*}{{\FR 17:13 }
Xeol é o lugar dos mortos
} será minha casa; nas trevas estenderei minha cama.
\VS{14}À cova chamo: Tu és meu pai; e aos vermes:
{\ADD{Sois}} minha mãe e minha irmã.
\VS{15}Onde, pois, estaria agora minha esperança? Quanto à minha esperança, quem a poderá ver?
\VS{16}Será que ela descerá aos ferrolhos do Xeol? Descansaremos juntos no pó da terra?

\par }\Chap{18}{\PP \VerseOne{1}Então Bildade, o suíta, respondeu, dizendo:
\VS{2}Quando é que dareis fim às palavras? Prestai atenção, e então falaremos.
\VS{3}Por que somos considerados animais, e tolos em vossos olhos?
\VS{4}Ó tu, que despedaças tua alma com tua ira; será a terra abandonada por tua causa, e será movida a rocha de seu lugar?
\VS{5}Na verdade, a luz dos perverso se apagará, e a faísca de seu fogo não brilhará.
\VS{6}A luz se escurecerá em sua tenda, e sua lâmpada sobre ele se apagará.
\VS{7}Os seus passos fortes serão encurtados, e seu
{\ADD{próprio}} intento o derrubará.
\VS{8}Pois será lançado à rede pelos seus
{\ADD{próprios}} pés, e sobre fios enredados andará.
\VS{9}Laço o pegará pelo calcanhar; a armadilha o prenderá.
\VS{10}Uma corda lhe está escondida debaixo da terra, e uma armadilha para ele está no caminho.
\VS{11}Assombros o espantarão ao redor, e o farão correr por onde seus passos forem.
\VS{12}Sua força se tornará em fome, e a perdição está pronta ao seu lado.
\FTNT{*}{{\FR 18:12 }
Sua força de tornará em fome
obscuro – trad. alt. Sua calamidade tem fome [dele]
}
\VS{13}Partes de sua pele serão consumidas; o primogênito da morte devorará os membros de seu corpo.
\VS{14}Será arrancado de sua tenda em que confiava, e será levado ao rei dos assombros.
\VS{15}Em sua tenda morará o que não é seu; enxofre se espalhará sobre a sua morada.
\VS{16}Por debaixo suas raízes se secarão, e por cima seus ramos serão cortados.
\VS{17}Sua memória perecerá da terra, e não terá nome pelas ruas.
\VS{18}Será lançado da luz para as trevas, e expulso será do mundo.
\VS{19}Não terá filho nem neto entre seu povo, nem sobrevivente em suas moradas.
\VS{20}Os do ocidente se espantarão com o seu dia, e os do oriente ficarão horrorizados.
\FTNT{†}{{\FR 18:20 }
os do ocidente, os do oriente
lit. os que vêm antes, os que vêm depois
}
\VS{21}Assim são as moradas do perverso, e este é o lugar
{\ADD{daquele que}} não reconhece a Deus.

\par }\Chap{19}{\PP \VerseOne{1}Porém Jó respondeu dizendo:
\VS{2}Até quando atormentareis minha alma, e me quebrantareis com palavras?
\VS{3}Já dez vezes me humilhastes; não tendes vergonha em me maltratar.
\VS{4}Mesmo se eu tiver errado, meu erro cabe apenas a mim.
\VS{5}Visto que vos exaltais contra mim, e contra mim usais minha desgraça,
\VS{6}Sabei, pois, que foi Deus que me transtornou, e
{\ADD{com}} sua rede me cercou.
\VS{7}Eis que eu clamo: Violência! Porém não sou respondido; grito, porém não há justiça.
\VS{8}Ele entrincheirou meu caminho, de modo que não consigo passar; e pôs trevas sobre minhas veredas.
\VS{9}Ele me despojou de minha honra, e tirou a coroa de minha cabeça.
\VS{10}Ele me derrubou por todos os lados, e pereço;
\FTNT{*}{{\FR 19:10 }
pereço
lit. e me vou
} e arrancou minha esperança como a uma árvore.
\VS{11}E fez inflamar contra mim sua ira, e me considerou para consigo como a
{\ADD{um de}} seus inimigos.
\VS{12}Juntas vieram suas tropas; prepararam contra mim seu caminho, e se acamparam ao redor de minha tenda.
\VS{13}Ele afastou meus irmãos para longe de mim; e os que me conheciam agora me estranham.
\VS{14}Meus parentes
{\ADD{me}} deixaram, e meus conhecidos se esqueceram de mim.
\VS{15}Os moradores de minha casa e minhas servas me tiveram por estranho; estrangeiro me tornei em seus olhos.
\VS{16}Chamei a meu servo, e ele não respondeu; de minha própria boca eu lhe suplicava.
\VS{17}Meu hálito é estranho à minha mulher, e sou repugnante aos filhos de minha mãe.
\VS{18}Até os meninos me desprezam; quando eu me levanto, falam contra mim.
\VS{19}Todos os meus amigos próximos me abominam; e
{\ADD{até}} aqueles que eu amava se viraram contra mim.
\VS{20}Meus ossos se grudaram à minha pele e à minha carne; e escapei
{\ADD{só}} com a pele de meus dentes.
\VS{21}Compadecei-vos de mim, meus amigos, compadecei-vos de mim; pois a mão de Deus me tocou.
\VS{22}Por que vós me perseguis como Deus, e não vos fartais de minhas carne?
\VS{23}Ah se minhas palavras fossem escritas! Ah se fossem escritas em um livro!
\VS{24}Que com ponta de ferro e com chumbo fossem esculpidas em pedra para sempre!
\VS{25}Pois eu sei que meu Redentor vive, e ao fim se levantará sobre a terra;
\VS{26}E mesmo depois de consumida minha pele, então em minha carne verei a Deus;
\VS{27}Ao qual eu verei para mim, e meus olhos
{\ADD{o}} verão, e não outro.
{\ADD{Isto é o que}} minhas entranhas
\FTNT{†}{{\FR 19:27 }
entranhas
lit. rins
} anseiam dentro de mim.
\VS{28}Se disserdes: Como o perseguiremos? Pois a raiz do problema se acha em mim,
\VS{29}Temei vós mesmos a espada; pois furor
{\ADD{há nos}} castigos pela espada; para que
{\ADD{assim}} saibais que
{\ADD{haverá}} julgamento.

\par }\Chap{20}{\PP \VerseOne{1}E Zofar, o naamita, respondeu, dizendo:
\VS{2}Por isso meus meus pensamentos me fazem responder; por causa da agitação dentro de mim.
\VS{3}Eu ouvi a repreensão que me envergonha; mas o espírito desde o meu entendimento responderá por mim.
\VS{4}Por acaso não sabes isto,
{\ADD{que foi}} desde a antiguidade, desde que o ser humano foi posto no mundo?
\VS{5}Que o júbilo dos perversos é breve, e a alegria do hipócrita
{\ADD{dura apenas}} um momento?
\VS{6}Ainda que sua altura subisse até o céu, e sua cabeça chegasse até as nuvens,
\VS{7}{\ADD{Mesmo assim}} com o seu excremento perecerá para sempre; os que houverem o visto, dirão: Onde ele está?
\VS{8}Como um sonho voará, e não será achado; e será afugentado como a visão noturna.
\VS{9}O olho que já o viu nunca mais o verá; nem seu lugar olhará mais para ele.
\VS{10}Seus filhos procurarão o favor dos pobres; e suas mãos devolverão a sua riqueza.
\VS{11}Seus ossos estão cheios de sua juventude, que juntamente com ele se deitará no pó.
\VS{12}Se o mal é doce em sua boca, e o esconde debaixo de sua língua;
\VS{13}Se o guarda para si, e não o abandona; ao contrário, o retém em sua boca.
\VS{14}Sua comida se mudará em suas entranhas, veneno de cobras será em seu interior.
\VS{15}Engoliu riquezas, porém as vomitará; Deus as tirará de seu ventre.
\VS{16}Veneno de cobras subará; língua de víbora o matará.
\VS{17}Não verá correntes, rios,
{\ADD{e}} ribeiros de mel e de manteiga.
\VS{18}Restituirá o trabalho e não o engolirá; da riqueza de seu comério não desfrutará.
\VS{19}Pois oprimiu
{\ADD{e}} desamparou aos pobres; roubou a casa que não edificou;
\VS{20}Por não ter sentido sossego em seu ventre, nada preservará de sua tão desejada riqueza.
\VS{21}Nada
{\ADD{lhe}} restou para que devorasse; por isso sua riqueza não será duradoura.
\VS{22}Estando cheio de sua fartura,
{\ADD{ainda}} estará angustiado; todo o poder
\FTNT{*}{{\FR 20:22 }
poder
lit. mão
} da miséria virá sobre ele.
\VS{23}Quando ele estiver enchendo seu vendre, Deus mandará sobre ele o ardor de sua ira, e
{\ADD{a}} choverá sobre ele em sua comida.
\VS{24}{\ADD{Ainda que}} fuja das armas de ferro, o arco de bronze o atravessará.
\VS{25}Ele a tirará de
{\ADD{seu}} corpo, e a ponta brilhante atingirá seu fígado; haverá sobre ele assombros.
\VS{26}Todas as trevas estão reservadas para seus tesouros escondidos; um fogo não assoprado o consumirá; acabará com o que restar em sua tenda.
\VS{27}Os céus revelarão sua maldade, e a terra se levantará contra ele.
\VS{28}As riquezas de sua casa serão transportadas; nos dias de sua ira elas se derramarão.
\VS{29}Esta é a parte que Deus dá ao homem perverso, a herança que Deus lhe prepara.

\par }\Chap{21}{\PP \VerseOne{1}Porém Jó respondeu, dizendo:
\VS{2}Ouvi atentamente minhas palavras, e seja isto vossas consolações.
\VS{3}Suportai-me, e eu falarei; e depois de eu ter falado,
{\ADD{então}} zombai.
\VS{4}Por acaso eu me queixo de algum ser humano? Porém ainda que
{\ADD{assim fosse}} ,por que meu espírito não se angustiaria?
\VS{5}Olhai-me, e espantai-vos; e ponde a mão sobre a boca.
\VS{6}Pois quando eu me lembro
{\ADD{disto}} ,me assombro, e minha carne é tomada de tremor.
\VS{7}Por que razão os perversos vivem, envelhecem, e ainda crescem em poder?
\VS{8}Seus filhos progridem com eles diante de seus rostos; e seus descendentes diante de seus olhos.
\VS{9}Suas casas têm paz, sem temor, e a vara de Deus não está contra eles.
\VS{10}Seus touros procriam, e não falham; suas vacas geram filhotes, e não abortam.
\VS{11}Suas crianças saem como um rebanho, e seus filhos saem dançando.
\VS{12}Levantam
{\ADD{a voz}} ao
{\ADD{som}} de tamboril e de harpa e se alegram ao som de flauta.
\VS{13}Em prosperidade gastam seus dias, e em um momento
\FTNT{*}{{\FR 21:13 }
em um momento
trad. alt. em paz
} descem ao Xeol.
\FTNT{†}{{\FR 21:13 }
Xeol é o lugar dos mortos
}
\VS{14}Assim dizem a Deus: Afasta-te de nós, porque não queremos conhecer teus caminhos.
\VS{15}Quem é o Todo-Poderoso, para que o sirvamos? E de que nos aproveitará que oremos a ele?
\VS{16}Eis que sua prosperidade não se deve às mãos deles. Longe de mim esteja o conselho dos perversos!
\VS{17}Quantas vezes sucede que a lâmpada dos perversos se apaga, e sua perdição vem sobre eles,
{\ADD{e}} Deus em sua ira
{\ADD{lhes}} reparte dores?
\VS{18}Eles serão como palha diante do vento, como o restos de palha que o turbilhão arrebata.
\VS{19}{\ADD{Vós dizeis}} : Deus guarda sua violência para seus filhos. Que
{\ADD{Deus}} pague ao próprio
{\ADD{perverso}} ,para que o conheça.
\VS{20}Seus olhos vejam sua ruína, e beba da ira do Todo-Poderoso.
\VS{21}Pois que interesse teria ele em sua casa depois de si, quando o número for cortado o número de seus meses?
\FTNT{‡}{{\FR 21:21 }
quando for cortado o número de seus meses
i.e., quando morrer
}
\VS{22}Poderia alguém ensinar conhecimento a Deus, que julga
{\ADD{até}} os que estão no alto?
\VS{23}Alguém morre na sua força plena, estando todo tranquilo e próspero,
\VS{24}Seus baldes estando cheios de leite, e o tutano de seus ossos umedecido.
\FTNT{§}{{\FR 21:24 }
baldes estão cheios de leite
obscuro – provável significado: seu corpo estando bem nutrido
}
\VS{25}Porém outro morre com amargura de alma, nunca tendo experimentado a prosperidade.
\VS{26}Juntamente jazem no pó, e os vermes os cobrem.
\VS{27}Eis que eu sei vossos pensamentos, e os mais intentos que planejais contra mim.
\VS{28}Porque dizeis: Onde está a casa do príncipe?, e: Onde está tenda das moradas dos perversos?
\VS{29}Por acaso não perguntastes aos que passam pelo caminho, e não conheceis seus sinais?
\VS{30}Que os maus são preservados no dia da destruição,
{\ADD{e}} são livrados no dia das fúrias?
\VS{31}Quem lhe denunciará seu caminho em sua face? E quem lhe pagará pelo que ele fez?
\VS{32}Finalmente ele é levado à sepultura, e no túmulo fazem vigilância.
\VS{33}Os torrões do vale lhe são doces; e todos o seguem; e adiante dele estão inúmeros.
\VS{34}Como, pois, me consolais em vão, já que vossas em vossas respostas
{\ADD{só}} resta falsidade?

\par }\Chap{22}{\PP \VerseOne{1}Então Elifaz, o temanita, respondeu, dizendo:
\VS{2}Por acaso o homem será de
{\ADD{algum}} proveito a Deus? Pode ele se beneficiar de algum sábio?
\VS{3}É útil ao Todo-Poderoso que sejas justo? Ganha ele algo se os teus caminhos forem íntegros?
\VS{4}Acaso ele te repreende
{\ADD{e}} vem contigo a juízo por causa da tua reverência?
\VS{5}{\ADD{Ou}} não será por causa de tua grande malícia, e de tuas maldades que não têm fim?
\VS{6}Porque tomaste penhor a teus irmãos sem causa, e foste tu que tiraste as roupas dos nus.
\VS{7}Não deste de beber água ao cansado, e negaste
\FTNT{*}{{\FR 22:7 }
negaste
lit. retiveste [contigo]
} o pão ao faminto.
\VS{8}Porém o homem poderoso teve a terra; e o homem influente habitava nela.
\VS{9}Às viúvas despediste vazias, e os braços dos órfãos foram quebrados.
\VS{10}Por isso que há laços ao redor de ti, e espanto repentino te perturbou;
\VS{11}Ou trevas, para que não vejas; e inundação de água te cobre.
\VS{12}Por acaso Deus não está na altura dos céus? Olha, pois, para o cume das estrelas, como estão elevadas.
\VS{13}Porém tu dizes: O que Deus sabe? Como ele julgará por entre a escuridão?
\VS{14}As nuvens são seu esconderijo, e ele não vê; ele passeia pela abóbada do céu.
\VS{15}Por acaso deste atenção para o velho caminho que pisaram os homens injustos?
\VS{16}Tais foram cortados antes de tempo;
{\ADD{sobre}} o fundamento deles foi derramada uma enchente.
\VS{17}Eles diziam a Deus: Afasta-te de nós! O que o Todo-Poderoso pode fazer por nós?
\VS{18}Sendo ele o que havia enchido suas casas de bens. Seja, porém, longe de mim o conselho dos perversos.
\VS{19}Os justos virão e se alegrarão; e o inocente os escarnecerá,
\VS{20}{\ADD{Dizendo}} : Certamente nossos adversários foram destruídos, e o que sobrou deles o fogo consumiu.
\VS{21}Reconcilia-te, pois, com
{\ADD{Deus}} ,e terás paz; assim o bem virá a ti.
\VS{22}Aceita, pois, a instrução
\FTNT{†}{{\FR 22:22 }
instrução
equiv. lei
} de sua boca, e põe suas palavras em teu coração.
\VS{23}Se te converteres ao Todo-Poderoso, serás edificado; se afastares a maldade de tua tenda,
\VS{24}E lançares
{\ADD{teu}} ouro no pó, o
{\ADD{ouro}} de Ofir junto às rochas dos ribeiros,
\VS{25}Então o próprio Todo-Poderoso será teu ouro, e tua prata maciça.
\VS{26}Porque então te deleitarás no Todo-Poderoso, e levantarás teu rosto a Deus.
\VS{27}Orarás a ele, e ele te ouvirá; e tu
{\ADD{lhe}} pagarás teus votos.
\VS{28}Aquilo que tu determinares se cumprirá a ti, e em teus caminhos a luz brilhará.
\VS{29}Quando
{\ADD{alguém}} for abatido, e tu disseres: Haja exaltação, Então
{\ADD{Deus}} salvará ao humilde.
\FTNT{‡}{{\FR 22:29 }
humilde
lit. abaixado de olhos
}
\VS{30}Ele libertará até ao que não é inocente, que será livrado pela pureza de tuas mãos.

\par }\Chap{23}{\PP \VerseOne{1}Porém Jó respondeu, dizendo:
\VS{2}Até hoje minha queixa é uma amargura; a mão
{\ADD{de Deus}} sobre mim é mais pesada que meu gemido.
\VS{3}Ah se eu soubesse como poderia achá-lo!
{\ADD{Então}} eu me chegaria até seu trono.
\VS{4}Apresentaria minha causa diante dele, e encheria minha boca de argumentos.
\VS{5}Eu saberia as palavras que ele me responderia, e entenderia o que me diria.
\VS{6}Por acaso ele brigaria comigo com seu grande poder? Não, pelo contrário, ele me daria atenção.
\VS{7}Ali o íntegro pleitearia com ele, e eu me livraria para sempre de meu Juiz.
\VS{8}Eis que se eu for ao oriente, ele não está ali;
{\ADD{se for}} ao ocidente, e não o percebo;
\VS{9}Se ao norte ele opera, eu não
{\ADD{o}} vejo; se ele se esconde ao sul, não
{\ADD{o}} enxergo.
\VS{10}Porém ele conhece meu caminho: Provar-me-á, e sairei como ouro.
\VS{11}Meus pés seguiram seus passos; guardei seu caminho, e não me desviei.
\VS{12}Nunca retirei
{\ADD{de mim}} o preceito de seus lábios, e guardei as palavras de sua boca mais que minha porção
{\ADD{de comida}} .
\VS{13}Porém se ele está decidido, quem poderá o desviar? O que sua alma quiser, isso fará.
\VS{14}Pois ele cumprirá o que está determinado para mim; ele
{\ADD{ainda}} tem muitas coisas como estas consigo.
\VS{15}Por isso eu me perturbo em sua presença. Quando considero
{\ADD{isto}} ,tenho medo dele.
\VS{16}Deus enfraqueceu meu coração; o Todo-Poderoso tem me perturbado.
\VS{17}Pois não estou destruído por causa das trevas, nem por causa da escuridão que encobriu meu rosto.

\par }\Chap{24}{\PP \VerseOne{1}Por que os tempos não são marcados pelo Todo-Poderoso? Por que os que o conhecem não veem seus dias?
\VS{2}Há os que mudam os limites de lugar, roubam rebanhos, e os apascentam.
\VS{3}Levam o asno do órfão; penhoram o boi da viúva.
\VS{4}Desviam do caminho aos necessitados; os pobres da terra juntos se escondem.
\VS{5}Eis que como asnos selvagens no deserto eles saem a seu trabalho buscando insistentemente por comida; o deserto dá alimento a ele
{\ADD{e a seus}} filhos.
\VS{6}No campo colhem sua forragem, e vindimam a vinha do perverso.
\VS{7}Passam a noite nus, por falta de roupa; sem terem coberta contra o frio.
\VS{8}Pelas correntes das montanhas são molhados e, não tendo abrigo, abraçam-se às rochas.
\VS{9}{\ADD{Há os que}} arrancam ao órfão do peito, e do pobre tomam penhor.
\VS{10}Ao nus fazem andar sem vestes, e fazem os famintos carregarem feixes.
\VS{11}Entre suas paredes espremem o azeite; pisam nas prensas de uvas, e
{\ADD{ainda}} têm sede.
\VS{12}Desde a cidade as pessoas gemem, e as almas dos feridos clamam; Mas Deus não dá atenção ao erro.
\VS{13}Há os que se opõem à luz; não conhecem seus caminhos, nem permanecem em suas veredas.
\VS{14}De manhã o homicida se levanta, mata ao pobre e ao necessitado, e de noite ele age como ladrão.
\VS{15}O olho do adúltero aguarda o crepúsculo, dizendo: Olho nenhum me verá; E esconde seu rosto.
\VS{16}Nas trevas vasculham as casas, de dia eles se trancam; não conhecem a luz.
\VS{17}Porque a manhã é para todos eles como sombra de morte; pois são conhecidos dos pavores de sombra de morte.
\VS{18}Ele é ligeiro sobre a superfície das águas; maldita é sua porção sobre a terra; não se vira para o caminho das vinhas.
\VS{19}A seca e o calor desfazem as águas da neve; assim
{\ADD{faz}} o Xeol
\FTNT{*}{{\FR 24:19 }
Xeol é o lugar dos mortos
} aos que pecaram.
\VS{20}A mãe se esquecerá dele; doce será para os vermes; nunca mais haverá memória
{\ADD{dele}} , e a perversidade será quebrada como um árvore.
\VS{21}Aflige à mulher estéril,
{\ADD{que}} não dá à luz; e nenhum bem faz à viúva.
\VS{22}Mas
{\ADD{Deus}} arranca aos poderosos com seu poder;
{\ADD{quando}} Deus se levanta, não há vida segura.
\VS{23}Se ele lhes dá descanso, nisso confiam;
{\ADD{mas}} os olhos de
{\ADD{Deus}} estão
{\ADD{postos}} nos caminhos deles.
\VS{24}São exaltados por um pouco
{\ADD{de tempo}} ,mas
{\ADD{logo}} desaparecem; são abatidos, encerrados como todos, e cortados como cabeças das espigas.
\VS{25}Se não é assim, quem me desmentirá, ou anulará minhas palavras?

\par }\Chap{25}{\PP \VerseOne{1}Então Bildade, o suíta, respondeu, dizendo:
\VS{2}O domínio e o temor estão com ele; ele faz paz em suas alturas.
\VS{3}Por acaso suas tropas têm número? E sobre quem não se levanta sua luz?
\VS{4}Como, pois, o ser humano seria justo para com Deus? E como seria puro aquele que nasce de mulher?
\VS{5}Eis que até a luz não tem brilho; nem as estrelas são puras diante de seus olhos.
\VS{6}Quanto menos o ser humano, que é uma larva, e o filho de homem, que é um verme.

\par }\Chap{26}{\PP \VerseOne{1}Porém Jó respondeu, dizendo:
\VS{2}Como tende ajudado ao que não tem força,
{\ADD{e}} sustentado ao braço sem vigor!
\VS{3}Como tende aconselhado ao que não tem conhecimento, e
{\ADD{lhe}} explicaste detalhadamente a verdadeira causa!
\VS{4}A quem tens dito
{\ADD{tais}} palavras? E de quem é o espírito que sai de ti?
\VS{5}Os mortos tremem debaixo das águas com os seus moradores.
\VS{6}O Xeol
\FTNT{*}{{\FR 26:6 }
Xeol é o lugar dos mortos
} está nu perante Deus
\FTNT{†}{{\FR 26:6 }
Lit. perante dele
} , e não há cobertura para a perdição.
\VS{7}Ele estende o norte sobre o vazio, suspende a terra sobre o nada.
\VS{8}Ele amarra as águas em suas nuvens, todavia a nuvem não se rasga debaixo dela.
\VS{9}Ele encobre a face de seu trono, e sobre ele estende sua nuvem.
\VS{10}Ele determinou limite à superfície das águas, até a fronteira entre a luz e as trevas.
\VS{11}As colunas do céu tremem, e se espantam por sua repreensão.
\VS{12}Ele agita o mar com seu poder, e com seu entendimento fere abate a Raabe.
\VS{13}Por seu Espírito
\FTNT{‡}{{\FR 26:13 }
Espírito
trad. alt. sopro
} adornou os céus; sua mão perfurou a serpente veloz.
\FTNT{§}{{\FR 26:13 }
veloz
trad. alt. fugitiva
}
\VS{14}Eis que estas são
{\ADD{somente}} as bordas de seus caminhos; e quão pouco é o que temos ouvido dele! Quem, pois, entenderia o trovão de seu poder?

\par }\Chap{27}{\PP \VerseOne{1}E Jó prosseguiu em falar seu discurso, e disse:
\VS{2}Vive Deus, que tirou meu direito, o Todo-Poderoso, que amargou minha alma,
\VS{3}Que enquanto meu fôlego estiver em mim, e o sopro de Deus em minhas narinas,
\VS{4}Meus lábios não falarão injustiça, nem minha língua pronunciará engano.
\VS{5}Nunca aconteça que eu diga que vós estais certos; até eu morrer
\FTNT{*}{{\FR 27:5 }
morrer
lit. entregue o espírito
} nunca tirarei de mim minha integridade.
\FTNT{†}{{\FR 27:5 }
nunca aconteça
lit. longe de mim
}
\VS{6}Eu me apegarei à minha justiça, e não a deixarei ir; meu coração não terá de que me acusar enquanto eu viver.
\FTNT{‡}{{\FR 27:6 }
enquanto eu viver
por [todos] os meus dias
}
\VS{7}Seja meu inimigo como o perverso, e o que se levantar contra mim como o injusto.
\VS{8}Pois qual é a esperança do hipócrita quando ele for cortado, quando Deus arrancar sua alma?
\VS{9}Por acaso Deus ouvirá seu clamor quando a aflição vier sobre ele?
\VS{10}Ele se deleitará no Todo-Poderoso? Invocará a Deus a todo tempo?
\VS{11}Eu vos ensinarei acerca da mão de Deus; não esconderei o que há com o Todo-Poderoso.
\VS{12}Eis que todos vós tendes visto
{\ADD{isso}} ; então por que vos deixais enganar por ilusão?
\VS{13}Esta é a porção do homem perverso para com Deus, a herança que os violentos receberão do Todo-Poderoso:
\VS{14}Se seus filhos se multiplicarem, serão para a espada; e seus descendentes não se fartarão de pão;
\VS{15}Os que lhe restarem, pela praga serão sepultados; e suas viúvas não chorarão.
\VS{16}Se ele amontoar prata como o pó da terra, e se preparar roupas como lama,
\VS{17}Mesmo ele tendo preparado, é o justo que se vestirá, e o inocente repartirá a prata.
\VS{18}Ele constrói sua casa como a traça, como uma barraca feita por um vigilante.
\VS{19}O rico dormirá, mas não será recolhido; ele abrirá seus olhos, e nada mais há para si.
\FTNT{§}{{\FR 27:19 }
nada mais há para si – obscuro
trad. alt. ele não existirá mais
}
\VS{20}Medos o tomarão como águas; um turbilhão o arrebatará de noite.
\VS{21}O vento oriental o levará, e ele partirá; e toma-o de seu lugar.
\VS{22}E o atacará sem o poupar,
{\ADD{enquanto}} ele tenta fugir de seu poder.
\FTNT{*}{{\FR 27:22 }
seu poder
lit. sua mão
}
\VS{23}Baterá palmas por causa dele, e desde seu lugar lhe assoviará.

\par }\Chap{28}{\PP \VerseOne{1}Certamente há minas para a prata, e o ouro lugar onde o derretem.
\VS{2}O ferro é tirado do solo, e da pedra se funde o cobre.
\VS{3}O
{\ADD{homem}} põe fim às trevas, e investiga em toda extremidade, as pedras que estão na escuridão e nas mais sombrias trevas.
\FTNT{*}{{\FR 28:3 }
mais sombrias trevas
lit. sombra de morte
}
\VS{4}Abre um poço onde não há morador, lugares esquecidos por quem passa a pé; pendurados longe da humanidade, vão de um lado para o outro.
\VS{5}Da terra o pão procede, e por debaixo ela é transformada como que pelo fogo.
\VS{6}Suas pedras são o lugar da safira, e contém pó de ouro.
\VS{7}A ave de rapina não conhece essa vereda; os olhos do falcão
\FTNT{†}{{\FR 28:7 }
falcão
obscuro – trad. alt. abutre
} não a viram.
\VS{8}Os filhotes de animais ferozes nunca a pisaram, nem o feroz leão passou por ela.
\VS{9}{\ADD{O homem}} põe sua mão no rochedo, e revolve os montes desde a raiz.
\VS{10}Cortou canais pelas rochas, e seus olhos veem tudo o que é precioso.
\VS{11}Ele tapa os rios desde suas nascentes, e faz o oculto sair para a luz.
\VS{12}Porém onde se achará a sabedoria? E onde está o lugar da inteligência?
\VS{13}O ser humano não conhece o valor dela, nem ela é achada na terra dos viventes.
\VS{14}O abismo diz: Não está em mim; E o mar diz: Nem comigo.
\VS{15}Nem por ouro fino se pode comprar, nem se pesar em troca de prata.
\VS{16}Não pode ser avaliada com ouro de Ofir, nem com ônix precioso, nem com safira.
\VS{17}Não se pode comparar com ela o ouro, nem o cristal; nem se pode trocar por joia de ouro fino.
\VS{18}De coral nem de quartzo não se fará menção; porque o preço da sabedoria é melhor que o de rubis.
\VS{19}O topázio de Cuxe não se pode comparar com ela; nem pode ser avaliada com o puro ouro fino.
\VS{20}De onde, pois, vem a sabedoria? E onde está o lugar da inteligência?
\VS{21}Porque encoberta está aos olhos de todo vivente, e é oculta a toda ave do céu.
\VS{22}O perdição e a morte dizem: Com nossos ouvidos ouvimos sua fama.
\VS{23}Deus entende o caminho dela, e ele conhece seu lugar.
\VS{24}Porque ele enxerga até os confins da terra, e vê tudo
{\ADD{o que há}} debaixo de céus.
\VS{25}Quando ele deu peso ao vento, e estabeleceu medida para as águas;
\VS{26}Quando ele fez lei para a chuva, e caminho para o relâmpago dos trovões,
\VS{27}Então ele a viu, e relatou; preparou-a, e também a examinou.
\VS{28}E disse ao homem: Eis que o temor ao Senhor é a sabedoria, e o desviar-se do mal
{\ADD{é}} a inteligência.

\par }\Chap{29}{\PP \VerseOne{1}E Jó continuou a falar seu discurso, dizendo:
\VS{2}Ah quem me dera que fosse como nos meses passados! Como nos dias em que Deus me guardava!
\VS{3}Quando ele fazia brilhar sua lâmpada sobre minha cabeça, e eu com sua luz caminhava pelas trevas,
\VS{4}Como era nos dias de minha juventude, quando a amizade de Deus estava sobre minha tenda;
\VS{5}Quando o Todo-Poderoso ainda estava comigo, meus filhos ao redor de mim;
\VS{6}Quando eu lavava meus passos com manteiga, e da rocha me corriam ribeiros de azeite!
\VS{7}Quando eu saía para a porta da cidade,
{\ADD{e}} na praça preparava minha cadeira,
\VS{8}Os rapazes me viam, e abriam caminho; e os idosos se levantavam, e ficavam em pé;
\VS{9}Os príncipes se detinham de falar, e punham a mão sobre a sua boca;
\VS{10}A voz dos líderes se calava, e suas línguas se apegavam a céu da boca;
\VS{11}O ouvido que me ouvia me considerava bem-aventurado, e o olho que me via dava bom testemunho de mim.
\VS{12}Porque eu livrava ao pobre que clamava, e ao órfão que não tinha quem o ajudasse.
\VS{13}A bênção do que estava a ponto de morrer vinha sobre mim; e eu fazia o coração da viúva ter grande alegria.
\VS{14}Vestia-me de justiça, e ela me envolvia; e meu juízo era como um manto e um turbante.
\VS{15}Eu era olhos para o cego, e pés para o manco.
\VS{16}Aos necessitados eu era pai; e a causa que eu não sabia, investigava com empenho.
\VS{17}E quebrava os queixos do perverso, e de seus dentes tirava a presa.
\VS{18}E eu dizia: Em meu ninho expirarei, e multiplicarei
{\ADD{meus}} dias como areia.
\VS{19}Minha raiz se estendia junto às águas, e o orvalho ficava de noite em meus ramos.
\VS{20}Minha honra se renovava em mim, e meu arco se revigorava em minha mão.
\VS{21}Ouviam-me, e esperavam; e se calavam ao meu conselho.
\VS{22}Depois de minha palavra nada replicavam, e minhas razões gotejavam sobre eles.
\VS{23}Pois esperavam por mim como pela chuva, e abriam sua boca como para a chuva tardia.
\VS{24}Se eu me ria com eles, não acreditavam; e não desfaziam a luz de meu rosto.
\VS{25}Eu escolhia o caminho para eles, e me sentava à cabeceira; e habitava como rei entre as tropas, como o consolador dos que choram.

\par }\Chap{30}{\PP \VerseOne{1}Porém agora riem de mim os mais jovens do que eu, cujos pais eu havia desdenhado até de os pôr com os cães de meu rebanho.
\VS{2}De que também me serviria força de suas mãos, nos quais o vigor já pereceu?
\VS{3}Por causa da pobreza e da fome andavam sós; roem na terra seca, no lugar desolado e deserto em trevas.
\VS{4}Que colhiam malvas entre os arbustos, e seu alimento eram as raízes dos zimbros.
\VS{5}Do meio
{\ADD{das pessoas}} eram expulsos, e gritavam contra eles, como a um ladrão.
\VS{6}Habitavam nos barrancos dos ribeiros secos, nos buracos da terra, e nas rochas.
\VS{7}Bramavam entre os arbustos, e se ajuntavam debaixo das urtigas.
\VS{8}Eram filhos de tolos, filhos sem nome, e expulsos de
{\ADD{sua}} terra.
\VS{9}Porém agora sirvo-lhes de chacota, e sou para eles um provérbio de escárnio.
\VS{10}Eles me abominam
{\ADD{e}} se afastam de mim; porém não hesitam em cuspir no meu rosto.
\VS{11}Pois
{\ADD{Deus}} desatou minha corda, e me oprimiu; por isso tiraram
{\ADD{de si}} todo constrangimento
\FTNT{*}{{\FR 30:11 }
constrangimento
lit. freio
} perante meu rosto.
\VS{12}À direita os jovens se levantam; empurram meus pés, e preparam contra mim seus caminhos de destruição.
\VS{13}Destroem meu caminho, e promovem minha miséria, sem necessitarem que alguém os ajude.
\VS{14}Eles vêm
{\ADD{contra mim}} como que por uma brecha larga,
{\ADD{e}} revolvem-se entre a desolação.
\VS{15}Pavores se voltam contra mim; perseguem minha honra como o vento, e como nuvem passou minha prosperidade.
\VS{16}Por isso agora minha alma se derrama em mim; dias de aflição têm me tomado.
\VS{17}De noite meus ossos se furam em mim, e meus pulsos não descansam.
\VS{18}Por grande força
{\ADD{de Deus}} minha roupa está estragada; ele me prendeu como a gola de minha roupa.
\VS{19}Lançou-me na lama, e fiquei semelhante ao pó e à cinza.
\VS{20}Clamo a ti, porém tu não me respondes; eu fico de pé, porém tu ficas
{\ADD{apenas}} olhando para mim.
\VS{21}Tu te tornaste cruel para comigo; com a força de tua mão tu me atacas.
\VS{22}Levantas-me sobre o vento,
{\ADD{e}} me fazes cavalgar
{\ADD{sobre ele}} ; e dissolves o meu ser.
\VS{23}Porque eu sei que me levarás à morte; e à casa determinada a todos os viventes.
\VS{24}Porém não se estende a mão para quem está em ruínas, quando clamam em sua opressão?
\VS{25}Por acaso eu não chorei pelo que estava em dificuldade,
{\ADD{e}} minha alma não se angustiou pelo necessitado?
\FTNT{†}{{\FR 30:25 }
que estava em dificuldades
lit. duro de dias
}
\VS{26}Quando eu esperava o bem, então veio o mal; quando eu esperava a luz, veio a escuridão.
\VS{27}Minhas entranhas fervem, e não se aquietam; dias de aflição me confrontam.
\VS{28}Ando escurecido, mas não pelo sol; levanto-me na congregação, e clamo por socorro.
\VS{29}Tornei-me irmão dos chacais, e companheiro dos avestruzes.
\FTNT{‡}{{\FR 30:29 }
avestruzes
obscuro – trad. alt. corujas
}
\VS{30}Minha pele se escureceu sobre mim, e meus ossos se inflamam de febre.
\VS{31}Por isso minha harpa passou a ser para lamentação, e minha flauta para vozes dos que choram.

\par }\Chap{31}{\PP \VerseOne{1}Eu fiz um pacto com meus olhos; como, pois, eu olharia com cobiça para a virgem?
\VS{2}Pois qual é a porção
{\ADD{dada}} por Deus acima, e a herança
{\ADD{dada}} pelo Todo-Poderoso das alturas?
\VS{3}Por acaso a calamidade não é para o perverso, e o desastre para os que praticam injustiça?
\VS{4}Por acaso ele não vê meus caminhos, e conta todos os meus passos?
\VS{5}Se eu andei com falsidade, e se meu pé se apressou para o engano,
\VS{6}Pese-me ele em balanças justas, e Deus saberá minha integridade.
\VS{7}Se meus passos se desviaram do caminho, e meu coração seguiu meus olhos, e se algo
\FTNT{*}{{\FR 31:7 }
algo
trad. alt. alguma impureza
} se apegou às minhas mãos,
\VS{8}Que eu semeie, e outro coma; e meus produtos sejam arrancados.
\VS{9}Se foi meu coração se deixou seduzir por
{\ADD{alguma}} mulher, ou se estive espreitei à porta de meu próximo,
\VS{10}Que minha mulher moa para outro, e outros se encurvem sobre ela.
\VS{11}Pois tal seria um crime vergonhoso, e delito
{\ADD{a ser sentenciado por}} juízes.
\VS{12}Pois seria um fogo que consumiria até à perdição, e destruiria toda a minha renda.
\VS{13}Se desprezei o direito de meu servo ou de minha serva quando eles reclamaram comigo,
\VS{14}Que faria eu quando Deus se levantasse? E quando ele investigasse
{\ADD{a causa}} ,o que eu lhe responderia?
\VS{15}Aquele que me fez no ventre
{\ADD{materno também}} não fez a ele? E não nos preparou de um mesmo
{\ADD{modo}} na madre?
\VS{16}Se eu neguei aos pobres o que eles desejavam, ou fiz desfalecer os olhos da viúva;
\VS{17}E se comi meu alimento sozinho, e o órfão não comeu dele
\VS{18}(Porque desde a minha juventude cresceu comigo como
{\ADD{se eu fosse seu pai}} ,e desde o ventre de minha mãe guiei
{\ADD{a viúva}} );
\VS{19}Se eu vi alguém morrer por falta de roupa, e o necessitado sem algo que o cobrisse,
\VS{20}Se sua cintura não me bendisse, quando ele se esquentava com as peles de meus cordeiros;
\VS{21}Se levantei minha mão contra o órfão, quando vi que seria favorecido na corte judicial,
\FTNT{†}{{\FR 31:21 }
seria favorecido na corte judicial – lit. teria ajuda na porta
a porta da cidade era onde os as causas judiciais eram julgadas
}
\VS{22}Que minha escápula caia do meu ombro, e meu braço se quebre de sua articulação.
\VS{23}Porque o castigo de Deus era um assombro para mim, e eu não teria poder contra sua majestade.
\VS{24}Se eu pus no ouro minha esperança, ou disse ao ouro fino: Tu és minha confiança;
\VS{25}Se eu me alegrei de que minha riqueza era muita, e de que minha mão havia obtido muito;
\VS{26}Se olhei para o sol quando brilhava, e à lua quando estava bela,
\VS{27}E meu coração se deixou enganar em segredo, e minha boca beijou minha mão,
\VS{28}Isto também seria um delito
{\ADD{a ser sentenciado por}} juiz; porque teria negado ao Deus de cima.
\VS{29}Se eu me alegrei da desgraça daquele que me odiava, e me agradei quando o mal o encontrou,
\VS{30}Sendo que nem deixei minha boca pecar, desejando sua morte com maldição,
\VS{31}Se a gente da minha casa nunca tivesse dito: Quem não se satisfez da carne dada por ele?
\VS{32}O estrangeiro não passava a noite na rua; eu abria minhas portas ao viajante.
\VS{33}Se encobri minhas transgressões como as pessoas
{\ADD{fazem}} ,
\FTNT{‡}{{\FR 31:33 }
as pessoas [fazem]
trad. alt. Adão [fez]
} escondendo meu delito em meu seio;
\VS{34}Porque eu tinha medo da grande multidão, e o desprezo das famílias me atemorizou; então me calei, e não saí da porta:
\VS{35}Quem me dera se alguém me ouvisse! Eis que minha vontade é que o Todo-Poderoso me responda, e meu adversário escrevesse um relato da acusação.
\VS{36}Certamente eu o carregaria sobre meu ombro, e o poria em mim como uma coroa.
\VS{37}Eu lhe diria o número de meus passos, e como um príncipe eu me chegaria a ele.
\VS{38}Se minha terra clamar contra mim, e seus sulcos juntamente chorarem;
\VS{39}Se comi seus frutos sem
{\ADD{pagar}} dinheiro, ou fiz expirar a alma de seus donos;
\VS{40}Em lugar de trigo que
{\ADD{me}} produza cardos, e ervas daninhas no lugar da cevada.
{\ADD{Aqui}} terminam as palavras de Jó.

\par }\Chap{32}{\PP \VerseOne{1}Então aqueles três homens cessaram de responder a Jó, porque ele era justo em seus olhos.
\VS{2}Porém se acendeu a ira de Eliú, filho de Baraquel, buzita, da família de Rão; contra Jó se acendeu sua ira, porque justificava mais a si mesmo que a Deus.
\VS{3}Também sua ira se acendeu contra seus três amigos, porque não achavam o que responder, ainda que tinham condenado a Jó.
\VS{4}E Eliú tinha esperado a Jó naquela discussão, porque tinham mais idade que ele.
\VS{5}Porém quando Eliú viu que não havia resposta na boca daqueles três homens, sua ira se acendeu.
\VS{6}Por isso Eliú, filho de Baraquel, buzita, respondeu, dizendo: Eu sou jovem e vós sois idosos; por isso fiquei receoso e tive medo de vos declarar minha opinião.
\VS{7}Eu dizia: Os dias falem, e a multidão de anos ensine sabedoria.
\VS{8}Certamente há espírito no ser humano, e a inspiração do Todo-Poderoso os faz entendedores.
\VS{9}Não são
{\ADD{somente}} os grandes que são sábios, nem
{\ADD{somente}} os velhos entendem o juízo.
\VS{10}Por isso eu digo: escutai-me; também eu declararei minha opinião.
\VS{11}Eis que eu aguardei vossas palavras,
{\ADD{e}} dei ouvidos a vossas considerações, enquanto vós buscáveis argumentos.
\VS{12}Eu prestei atenção a vós, porém eis que ninguém há de vós que possa convencer a Jó,
{\ADD{nem}} que responda a suas palavras.
\VS{13}Portanto não digais: Encontramos a sabedoria; que Deus o derrote,
{\ADD{e}} não o homem.
\VS{14}Jó não dirigiu suas palavras a mim, nem eu lhe responderei com vossos dizeres.
\VS{15}Estão pasmos, não respondem mais; faltam-lhes palavras.
\VS{16}Esperei, pois, porém
{\ADD{agora}} não falam; porque já pararam,
{\ADD{e}} não respondem mais.
\VS{17}Também eu responderei minha parte; também eu declararei minha opinião.
\VS{18}Porque estou cheio de palavras, e o espírito em meu ventre me obriga.
\VS{19}Eis que meu ventre é como o vinho que não tem abertura; e que está a ponto de arrebentar como odres novos.
\VS{20}Falarei, para que eu me alivie; abrirei meus lábios, e responderei.
\VS{21}Não farei eu acepção de pessoas, nem usarei de títulos lisonjeiros para com o homem.
\VS{22}Pois não sei lisonjear;
{\ADD{caso contrário}} ,meu Criador logo me removeria.

\par }\Chap{33}{\PP \VerseOne{1}Portanto, Jó, ouve, por favor, meus dizeres, e dá ouvidos a todas as minhas palavras.
\VS{2}Eis que já abri minha boca; minha língua já fala debaixo do meu céu da boca.
\VS{3}Meus dizeres pronunciarão a integridade do meu coração, e o puro conhecimento dos meus lábios.
\VS{4}O Espírito de Deus me fez, e o sopro do Todo-Poderoso me deu vida.
\VS{5}Se puderes, responde-me; dispõe-te perante mim, e persiste.
\VS{6}Eis que para Deus eu sou como tu; do barro também eu fui formado.
\VS{7}Eis que meu terror não te espantará, nem minha mão será pesada sobre ti.
\VS{8}Certamente tu disseste a meus ouvidos, e eu ouvi a voz de tuas palavras,
\VS{9}{\ADD{Que diziam}} : Eu sou limpo e sem transgressão; sou inocente, e não tenho culpa.
\VS{10}Eis que
{\ADD{Deus}} buscou pretextos contra mim,
{\ADD{e}} me tem por seu inimigo.
\VS{11}Ele pôs meus pés no tronco, e observa todas as minhas veredas.
\VS{12}Eis que nisto não foste justo, eu te respondo; pois Deus é maior que o ser humano.
\VS{13}Por que razão brigas contra ele por não dar resposta às palavras do ser humano?
\FTNT{*}{{\FR 33:13 }
do ser humano
lit. palavras dele
}
\VS{14}Contudo Deus fala uma ou duas vezes, ainda que
{\ADD{o ser humano}} não entenda.
\VS{15}Em sonho
{\ADD{ou em}} visão noturna, quando o sono profundo cai sobre as pessoas,
{\ADD{e}} adormecem na cama.
\VS{16}Então o revela ao ouvido das pessoas, e os sela com advertências;
\VS{17}Para desviar ao ser humano de sua obra, e do homem a soberba.
\VS{18}Para desviar a sua alma da perdição, e sua vida de passar pela espada.
\VS{19}Também em sua cama é castigado com dores, com luta constante em seus ossos,
\VS{20}De modo que sua vida detesta
{\ADD{até}} o pão, e sua alma a comida deliciosa.
\VS{21}Sua carne desaparece da vista, e seus ossos, que antes não se viam, aparecem.
\VS{22}Sua alma se aproxima da cova, e sua vida dos que causam a morte.
\VS{23}Se com ele, pois, houver algum anjo, algum intérprete; um dentre mil, para anunciar ao ser humano o que lhe é correto,
\VS{24}Então
{\ADD{Deus}} terá misericórdia dele, e
{\ADD{lhe}} dirá: Livra-o, para que não desça à perdição;
{\ADD{já}} achei o resgate.
\VS{25}Sua carne se rejuvenescerá mais do que era na infância,
{\ADD{e}} voltará aos dias de sua juventude.
\VS{26}Ele orará a Deus, que se agradará dele; e verá sua face com júbilo, porque ele restituirá ao ser humano sua justiça.
\VS{27}Ele olhará para as pessoas, e dirá: Pequei, e perverti o
{\ADD{que era}} correto, o que de nada me aproveitou.
\VS{28}{\ADD{Porém}} Deus livrou minha alma para que eu não passasse à cova, e
{\ADD{agora}} minha vida vê a luz!
\VS{29}Eis que Deus faz tudo isto duas
{\ADD{ou}} três vezes com o ser humano,
\VS{30}Para desviar sua alma da perdição, e o iluminar com a luz dos viventes.
\VS{31}Presta atenção, Jó, e ouve-me; cala-te, e eu falarei.
\VS{32}Se tiveres o que dizer, responde-me; fala, porque eu quero te justificar.
\VS{33}E se não, escuta-me; cala-te, e eu ensinarei sabedoria.

\par }\Chap{34}{\PP \VerseOne{1}Eliú respondeu mais, dizendo:
\VS{2}Ouvi, vós sábios, minhas palavras; e vós, inteligentes, dai-me ouvidos.
\VS{3}Porque o ouvido prova as palavras, assim como o paladar experimenta a comida.
\VS{4}Escolhamos para nós o que é correto,
{\ADD{e}} conheçamos entre nós o que é bom.
\VS{5}Pois Jó disse: Eu sou justo, e Deus tem me tirado meu direito.
\VS{6}Por acaso devo eu mentir quanto ao meu direito? Minha ferida
\FTNT{*}{{\FR 34:6 }
ferida
lit. flecha
} é dolorosa mesmo que eu não tenha transgressão.
\VS{7}Que homem há como Jó, que bebe o escárnio como água?
\VS{8}E que caminha na companhia dos que praticam maldade, e anda com homens perversos?
\VS{9}Porque disse: De nada aproveita ao homem agradar-se em Deus.
\FTNT{†}{{\FR 34:9 }
agradar-se em Deus
trad. alt. ser do agrado de Deus
}
\VS{10}Portanto vós, homens de bom-senso, escutai-me; longe de Deus esteja a maldade, e do Todo- Poderoso a perversidade!
\VS{11}Porque ele paga ao ser humano
{\ADD{conforme}} sua obra, e faz a cada um conforme o seu caminho.
\VS{12}Certamente Deus não faz injustiça, e o Todo-Poderoso não perverte o direito.
\VS{13}Quem o pôs para administrar a terra? E quem dispôs a todo o mundo?
\VS{14}Se ele tomasse a decisão,
\FTNT{‡}{{\FR 34:14 }
tomasse a decisão
lit. pusesse seu coração em si
} e recolhesse para si seu espírito e seu fôlego,
\VS{15}Toda carne juntamente expiraria, e o ser humano se tornaria em pó.
\VS{16}Se pois há
{\ADD{em ti}} entendimento, ouve isto: dá ouvidos ao som de minhas palavras.
\VS{17}Por acaso quem odeia a justiça poderá governar? E condenarás tu ao Poderoso Justo?
\VS{18}Pode, por acaso, o rei ser chamado de vil,
{\ADD{e}} os príncipes de perversos?
\VS{19}{\ADD{Quanto menos}} a aquele que não faz acepção de pessoas de príncipes, nem valoriza mais o rico que o pobre! Pois todos são obras de suas mãos.
\VS{20}Em um momento morrem; e à meia noite os povos são sacudidos, e passam; e o poderoso será tomado sem ação humana.
\FTNT{§}{{\FR 34:20 }
ação humana
lit. mão
}
\VS{21}Porque seus olhos estão sobre os caminhos do homem, e ele vê todos os seus passos.
\VS{22}Não há trevas nem sombra de morte em que os que praticam maldade possam se esconder.
\VS{23}Pois ele não precisa observar tanto ao homem para que este possa entrar em juízo com Deus.
\VS{24}Ele quebranta aos fortes sem
{\ADD{precisar de}} investigação, e põe outros em seu lugar.
\VS{25}Visto que ele conhece suas obras, de noite os trastorna, e ficam destroçados.
\VS{26}Ele os espanca à vista pública por serem maus.
\VS{27}Pois eles se desviaram de segui-lo, e não deram atenção a nenhum de seus caminhos.
\VS{28}Assim fizeram com que viesse a ele o clamor do pobre, e ele ouvisse o clamor dos aflitos.
\VS{29}E se ele ficar quieto, quem
{\ADD{o}} condenará? Se ele esconder o rosto, quem o olhará?
{\ADD{Ele está}} quer sobre um povo, quer sobre um único ser humano,
\VS{30}Para que a pessoa hipócrita não reine,
{\ADD{e}} não haja ciladas ao povo.
\VS{31}Por que
{\ADD{não tão somente}} se diz: Suportei
{\ADD{teu castigo}} não farei mais o que é errado;
\VS{32}Ensina-me o que não vejo; se fiz alguma maldade, nunca mais a farei?
\VS{33}Será a recompensa da parte dele como tu queres, para que a recuses? És tu que escolhes, e não eu; o que tu sabes, fala.
\VS{34}As pessoas de entendimento dirão comigo, e o homem sábio me ouvirá;
\VS{35}Jó não fala com conhecimento, e a suas palavras falta prudência.
\VS{36}Queria eu que Jó fosse provado até o fim, por causa de suas respostas comparáveis a de homens malignos.
\VS{37}Pois ao seu pecado ele acrescentou rebeldia; ele bate as mãos
{\ADD{de forma desrespeitosa}} entre nós, e multiplica suas palavras contra Deus.

\par }\Chap{35}{\PP \VerseOne{1}Eliú respondeu mais, dizendo:
\VS{2}Pensas tu ser direito dizeres: Maior é a minha justiça do que a de Deus?
\VS{3}Porque disseste: Para que ela te serve?
{\ADD{Ou}} : Que proveito terei dela mais que meu pecado?
\VS{4}Eu darei reposta a ti, e a teus amigos contigo.
\VS{5}Olha para os céus, e vê; e observa as nuvens,
{\ADD{que}} são mais altas que tu.
\VS{6}Se tu pecares, que
{\ADD{mal}} farás contra ele? Se tuas transgressões se multiplicarem, que
{\ADD{mal}} lhe farás?
\VS{7}Se fores justo, que lhe darás? Ou o que ele receberá de tua mão?
\VS{8}Tua perversidade
{\ADD{poderia afetar}} a outro homem como tu; e tua justiça
{\ADD{poderia ser proveitosa}} a
{\ADD{algum}} filho do homem.
\VS{9}{\ADD{Os aflitos}} clamam por causa da grande opressão; eles gritam por causa do poder
\FTNT{*}{{\FR 35:9 }
poder
lit. braço
} dos grandes.
\VS{10}Porém ninguém diz: Onde está Deus, meu Criador, que dá canções na noite,
\VS{11}Que nos ensina mais que aos animais da terra, e nos faz sábios mais que as aves do céu?
\VS{12}Ali clamam, porém ele não responde, por causa da arrogância dos maus.
\VS{13}Certamente Deus não ouvirá a súplica vazia, nem o Todo-Poderoso dará atenção a ela.
\VS{14}Quanto menos ao que disseste: que tu não o vês! Porém o juízo está diante dele; portanto espera nele.
\FTNT{†}{{\FR 35:14 }
v. 14
obscuro – trad. alt. Quanto menos ao que disseste: que tu não o vês, [que] o juízo está diante dele, e [que] tu nele esperas.
}
\VS{15}Mas agora, já que a ira dele
{\ADD{ainda}} não está castigando, e ele não deu completa atenção à arrogância,
\FTNT{‡}{{\FR 35:15 }
v. 15 – obscuro
trad. alt. Ainda mais: que a ira dele não castiga, e que ele não presta atenção à arrogância.
}
\VS{16}Por isso Jó abriu sua boca em vão, e multiplicou palavras sem conhecimento.

\par }\Chap{36}{\PP \VerseOne{1}Prosseguiu Eliú ainda, dizendo:
\VS{2}Espera-me um pouco, e eu te mostrarei que ainda há palavras a favor de Deus.
\VS{3}Desde longe trarei meu conhecimento, e a meu Criador atribuirei a justiça.
\VS{4}Porque verdadeiramente minhas palavras não serão falsas; contigo está um que tem completo conhecimento.
\VS{5}Eis que Deus é grande, porém despreza ninguém; grande ele é em poder de entendimento.
\FTNT{*}{{\FR 36:5 }
entendimento
lit. coração
}
\VS{6}Ele não permite o perverso viver, e faz justiça aos aflitos.
\VS{7}Ele não tira seus olhos do justo; ao contrário, ele os faz sentar com os reis no trono, e
{\ADD{assim}} são exaltados.
\VS{8}E se estiverem presos em grilhões, e detidos com cordas de aflição,
\VS{9}Então ele lhes faz saber as obras que fizeram, e suas transgressões, das quais se orgulharam.
\VS{10}E revela a seus ouvidos, para que sejam disciplinados; e lhes diz, para que se convertam da maldade.
\VS{11}Se ouvirem, e
{\ADD{o}} servirem, acabarão seus dias em prosperidade, e seus anos em prazeres.
\VS{12}Porém se não ouvirem, perecerão pela espada, e morrerão sem conhecimento.
\VS{13}E os hipócritas de coração acumulam a ira
{\ADD{divina}} ; e quando ele os amarrar, mesmo assim não clamam.
\VS{14}A alma deles morrerá em sua juventude, e sua vida entre os pervertidos.
\FTNT{†}{{\FR 36:14 }
pervertidos
provavelmente prostitutos que faziam relações sexuais em rituais idólatras
}
\VS{15}Ele livra o aflito de sua aflição, e na opressão ele revela a seus ouvidos.
\VS{16}Assim também ele pode te desviar da boca da angústia
{\ADD{para}} um lugar amplo, onde não haveria aperto; para o conforto de tua mesa, cheia dos melhores alimentos.
\FTNT{‡}{{\FR 36:16 }
melhores alimentos
lit. gordura
}
\VS{17}Mas tu estás cheio do julgamento do perverso; o julgamento e a justiça te tomam.
\VS{18}Por causa da furor,
{\ADD{guarda-te}} para que não sejas seduzido pela riqueza, nem que um grande suborno
\FTNT{§}{{\FR 36:18 }
suborno
lit. resgate
} te faça desviar.
\FTNT{*}{{\FR 36:18 }
[cuida] para que não sejas seduzido pela riqueza – obscuro
trad. alt. [Cuida] para que o furor não te induza a escarnecer
}
\VS{19}Pode, por acaso, a tua riqueza
\FTNT{†}{{\FR 36:19 }
a tua riqueza
obscuro – trad. alt. o teu clamor
} te sustentar para que não tenhas aflição, mesmo com todos os esforços de
{\ADD{teu}} poder?
\VS{20}Não anseies pela noite, em que os povos são tomados de seu lugar.
\VS{21}Guarda-te, e não te voltes para a maldade; pois por isto que tens sido testado com miséria.
\FTNT{‡}{{\FR 36:21 }
pois por isto que tens sido testado com miséria
trad. alt. pois preferiste isto à miséria
}
\VS{22}Eis que Deus é exaltado em seu poder; que instrutor há como ele?
\VS{23}Quem lhe indica o seu caminho? Quem poderá lhe dizer: Cometeste maldade?
\VS{24}Lembra-te de engrandeceres sua obra, a qual os seres humanos contemplam.
\VS{25}Todas as pessoas a veem; o ser humano a enxerga de longe.
\VS{26}Eis que Deus é grande, e nós não o compreendemos; não se pode descobrir o número de seus anos.
\VS{27}Ele traz para cima as gotas das águas, que derramam a chuva de seu vapor;
\VS{28}A qual as nuvens destilam, gotejando abundantemente sobre o ser humano.
\VS{29}Poderá alguém entender a extensão das nuvens,
{\ADD{e}} os estrondos de seu pavilhão?
\VS{30}Eis que estende sobre ele sua luz, e cobre as profundezas
\FTNT{§}{{\FR 36:30 }
profundezas
lit. raízes
} do mar.
\VS{31}Pois por estas coisas ele julga aos povos, e dá alimento em abundância.
\VS{32}Ele cobre as mãos com o relâmpago, e dá ordens para que atinja o alvo.
\VS{33}O trovão anuncia sua presença; o gado também
{\ADD{prenuncia a tempestade}} que se aproxima.

\par }\Chap{37}{\PP \VerseOne{1}Disto também o meu coração treme, e salta de seu lugar.
\VS{2}Ouvi atentamente o estrondo de sua voz, e o som que sai de sua boca,
\VS{3}Ao qual envia por debaixo de todos os céus; e sua luz até os confins da terra.
\VS{4}Depois disso brama com estrondo; troveja com sua majestosa voz; e ele não retém
{\ADD{seus relâmpagos}} quando sua voz é ouvida.
\VS{5}Deus troveja maravilhosamente com sua voz; ele faz coisas tão grandes que nós não compreendemos.
\VS{6}Pois ele diz à neve: Cai sobre à terra; Como também à chuva: Sê chuva forte.
\VS{7}Ele sela as mãos de todo ser humano, para que todas as pessoas conheçam sua obra.
\VS{8}E os animais selvagens entram nos esconderijos, e ficam em suas tocas.
\VS{9}Da recâmara
\FTNT{*}{{\FR 37:9 }
recâmara
trad. alt. sul
} vem o redemoinho, e dos
{\ADD{ventos}} que espalham
\FTNT{†}{{\FR 37:9 }
[ventos] que espalham
trad. alt. norte
}
{\ADD{vem}} o frio.
\VS{10}Pelo sopro de Deus se dá o gelo, e as largas águas se congelam.
\VS{11}Ele também carrega de umidade as espessas nuvens,
{\ADD{e}} por entre as nuvens ele espalha seu relâmpago.
\FTNT{‡}{{\FR 37:11 }
relâmpago
lit. luz – também no resto do capítulo
}
\VS{12}Então elas se movem ao redor segundo sua condução, para que façam quanto ele lhes manda sobre a superfície do mundo, na terra;
\VS{13}Seja que ou por vara de castigo, ou para sua terra, ou por bondade as faça vir.
\VS{14}Escuta isto, Jó; fica parado, e considera as maravilhas de Deus.
\VS{15}Por acaso sabes tu quando Deus dá ordem a elas, e faz brilhar o relâmpago de sua nuvem?
\VS{16}Conheces tu os equilíbrios das nuvens, as maravilhas daquele que é perfeito no conhecimento?
\VS{17}Tu, cujas vestes se aquecem quando a terra se aquieta por causa do
{\ADD{vento}} sul,
\VS{18}acaso podes estender com ele os céus, que estão firmes como um espelho fundido?
\VS{19}Ensina-nos o que devemos dizer a ele;
{\ADD{pois discurso}} nenhum podemos propor, por causa das
{\ADD{nossas}} trevas.
\VS{20}Seria contado a ele o que eu haveria de falar? Por acaso alguém falaria para ser devorado?
\FTNT{§}{{\FR 37:20 }
para ser devorado
obscuro – trad. alt. quando está confuso
}
\VS{21}E agora não se pode olhar para o sol, quando brilha nos céus, quando o vento passa e os limpa.
\VS{22}Do norte vem o esplendor dourado; em Deus há majestade temível.
\VS{23}Não podemos alcançar ao Todo-Poderoso; ele é grande em poder; porém ele a ninguém oprime
{\ADD{em}} juízo e grandeza de justiça.
\VS{24}Por isso as pessoas o temem; ele não dá atenção aos que
{\ADD{se acham}} sábios de coração.

\par }\Chap{38}{\PP \VerseOne{1}Então o SENHOR respondeu a Jó desde um redemoinho, e disse:
\VS{2}Quem é esse que obscurece o conselho com palavras sem conhecimento?
\VS{3}Agora cinge teus lombos como homem; e eu te perguntarei, e tu me explicarás.
\VS{4}Onde estavas tu quando eu fundava a terra? Declara-
{\ADD{me}} ,se tens inteligência.
\VS{5}Quem determinou suas medidas, se tu o sabes? Ou quem estendeu cordel sobre ela?
\VS{6}Sobre o que estão fundadas suas bases? Ou quem pôs sua pedra angular,
\VS{7}Quando as estrelas do amanhecer cantavam alegremente juntas, e todos os filhos de Deus jubilavam?
\VS{8}Ou
{\ADD{quem}} encerrou o mar com portas, quando transbordou, saindo da madre,
\VS{9}Quando eu pus nuvens por sua vestidura, e a escuridão por sua faixa;
\VS{10}Quando eu passei sobre ele meu decreto, e
{\ADD{lhe}} pus portas e ferrolhos,
\VS{11}E disse: Até aqui virás, e não passarás adiante, e aqui será o limite para a soberba de tuas ondas?
\VS{12}Desde os teus dias tens dado ordem à madrugada?
{\ADD{Ou}} mostraste tu ao amanhecer o seu lugar,
\VS{13}Para que tomasse os confins da terra, e os perversos fossem sacudidos dela?
\VS{14}E
{\ADD{a terra}} se transforma como barro
{\ADD{sob}} o selo;
{\ADD{todas as coisas sobre ela}} se apresentam como vestidos?
\VS{15}E dos perversos é desviada sua luz, e o braço erguido é quebrado.
\VS{16}Por acaso chegaste tu às fontes do mar, ou passeaste no mais profundo abismo?
\VS{17}Foram reveladas a ti as portas da morte, ou viste as portas da sombra de morte?
\VS{18}Entendeste tu as larguras da terra? Declara, se sabes tudo isto.
\VS{19}Onde está o caminho
{\ADD{por onde}} mora a luz? E quanto às trevas, onde fica o seu lugar?
\VS{20}Para que as tragas a seus limites, e conheças os caminhos de sua casa.
\VS{21}Certamente tu o sabes, pois já eras nascido, e teus dias são inúmeros!
\VS{22}Por acaso entraste tu aos depósitos
\FTNT{*}{{\FR 38:22 }
depósitos
lit. tesouros
} da neve, e viste os depósitos do granizo,
\VS{23}Que eu retenho até o tempo da angústia, até o dia da batalha e da guerra?
\VS{24}Onde está o caminho em que a luz
\FTNT{†}{{\FR 38:24 }
luz
trad. alt. relâmpago
} se reparte, e o vento oriental se dispersa sobre a terra?
\VS{25}Quem repartiu um canal às correntezas de águas, e caminho aos relâmpagos dos trovões,
\VS{26}Para chover sobre a terra
{\ADD{onde}} havia ninguém,
{\ADD{sobre}} o deserto, onde não há gente,
\VS{27}Para fartar
{\ADD{a terra}} deserta e desolada, e para fazer crescer aos renovos da erva.
\VS{28}Por acaso a chuva tem pai? Ou quem gera as gotas do orvalho?
\VS{29}De qual ventre procede o gelo? E quem gera a geada do céu?
\VS{30}As águas se tornam duras como pedra, e a superfície do abismo se congela.
\VS{31}Podes tu atar as cadeias das Plêiades, ou desatar as cordas do Órion?
\VS{32}Podes tu trazer as constelações
\FTNT{‡}{{\FR 38:32 }
constelações
obscuro – hebraico Mazarote
} a seu tempo, e guiar a Ursa com seus filhos?
\VS{33}Sabes tu as ordenanças dos céus? Ou podes tu dispor do domínio deles sobre a terra?
\VS{34}Ou podes levantar tua voz até às nuvens, para que a abundância das águas te cubra?
\VS{35}Podes tu mandar relâmpagos, para que saiam, e te digam: Eis-nos aqui?
\VS{36}Quem pôs a sabedoria no íntimo? Ou quem deu entendimento à mente?
\VS{37}Quem pode enumerar as nuvens com sabedoria? E os odres dos céus, quem pode os despejar?
\VS{38}Quando o pó se endurece, e os torrões se apegam uns aos outros?
\VS{39}Caçarás tu a presa para o leão? Ou saciarás a fome dos leões jovens,
\VS{40}Quando estão agachados nas covas,
{\ADD{ou}} estão à espreita no matagal?
\VS{41}Quem prepara aos corvos seu alimento, quando seus filhotes clamam a Deus, andando de um lado para o outro por não terem o que comer?

\par }\Chap{39}{\PP \VerseOne{1}Sabes tu o tempo em que as cabras montesas dão filhotes? Ou observaste tu as cervas quando em trabalho de parto?
\VS{2}Contaste os meses que elas cumprem, e sabes o tempo de seu parto?
\VS{3}Quando se encurvam, produzem seus filhos,
{\ADD{e}} lançam de si suas dores.
\VS{4}Seus filhos se fortalecem, crescem como o trigo; saem, e nunca mais voltam a elas.
\VS{5}Quem despediu livre ao asno montês? E quem ao asno selvagem soltou das ataduras?
\VS{6}Ao qual eu dei a terra desabitada por casa, e a terra salgada por suas moradas.
\VS{7}Ele zomba do tumulto da cidade; não ouve os gritos do condutor.
\VS{8}A extensão dos montes é seu pasto; e busca tudo o que é verde.
\VS{9}Por acaso o boi selvagem
\FTNT{*}{{\FR 39:9 }
boi selvagem
Almeida 1819: unicórnio
} quererá te servir, ou ficará junto de tua manjedoura?
\VS{10}Ou amarrarás ao boi selvagem com sua corda para o arado? Ou lavrará ele aos campos atrás de ti?
\VS{11}Confiarás nele, por ser grande sua força, e deixarás que ele faça teu trabalho?
\VS{12}Porás tua confiança nele, para que ele colha tua semente, e a junte em tua eira?
\VS{13}As asas da avestruz batem alegremente, mas são suas asas e penas como as da cegonha?
\VS{14}Ela deixa seus ovos na terra, e os esquenta no chão,
\VS{15}E se esquece de que pés podem os pisar, e os animais do campo
{\ADD{podem}} os esmagar.
\VS{16}Age duramente para com seus filhos, como se não fossem seus, sem temer que seu trabalho tenha sido em vão.
\VS{17}Porque Deus a privou de sabedoria, e não lhe repartiu entendimento.
\VS{18}Porém quando se levanta para correr,
\FTNT{†}{{\FR 39:18 }
levanta para correr – significado obscuro
lit. levanta no alto
} zomba do cavalo e do seu cavaleiro.
\VS{19}És tu que dás força ao cavalo, ou que vestes seu pescoço com crina?
\VS{20}Podes tu o espantar como a um gafanhoto? O sopro de suas narinas é terrível.
\VS{21}Ele escarva a terra, alegra-se de sua força,
{\ADD{e}} sai ao encontro das armas;
\VS{22}Ele zomba do medo, e não se espanta; nem volta para trás por causa da espada.
\VS{23}Contra ele rangem a aljava, o ferro brilhante da lança e do dardo;
\VS{24}Sacudido-se com furor, ele escarva a terra; ele não fica parado ao som da trombeta.
\VS{25}Ao som das trombetas diz: Eia! E desde longe cheira a batalha, o grito dos capitães, e o barulho.
\VS{26}Por acaso é por tua inteligência que o gavião voa,
{\ADD{e}} estende suas asas para o sul?
\VS{27}Ou é por tua ordem que a água voa alto e põe seu ninho na altura?
\VS{28}Nas penhas ela mora e habita; no cume das penhas, e em lugares seguros.
\VS{29}Desde ali espia a comida; seus olhos avistam de longe.
\VS{30}Seus filhotes sugam sangue; e onde houver cadáveres, ali ela está.

\par }\Chap{40}{\PP \VerseOne{1}Então o SENHOR respondeu mais a Jó, dizendo:
\VS{2}Por acaso quem briga contra o Todo-Poderoso pode ensiná-lo? Quem quer repreender a Deus, responda a isto.
\VS{3}Então Jó respondeu ao SENHOR, dizendo:
\VS{4}Eis que eu sou insignificante; o que eu te responderia? Ponho minha mão sobre minha boca.
\VS{5}Uma vez falei, porém não responderei; até duas vezes, porém não prosseguirei.
\VS{6}Então o SENHOR respondeu a Jó desde o redemoinho, dizendo:
\VS{7}Cinge-te agora os teus lombos como homem; eu te perguntarei, e tu me explica.
\VS{8}Por acaso tu anularias o meu juízo? Tu me condenarias, para te justificares?
\VS{9}Tens tu braço como Deus? Ou podes tu trovejar com
{\ADD{tua}} voz como ele?
\VS{10}Orna-te, pois, de excelência e alteza; e veste-te de majestade e glória.
\VS{11}Espalha os furores de tua ira; olha a todo soberbo, e abate-o.
\VS{12}Olha a todo soberbo, e humilha-o; e esmaga aos perversos em seu lugar.
\VS{13}Esconde-os juntamente no pó; ata seus rostos no oculto.
\VS{14}E eu também te reconhecerei; pois tua mão direita te terá livrado.
\VS{15}Observa o beemote,
\FTNT{*}{{\FR 40:15 }
beemote
obscuro
} ao qual eu fiz contigo; ele come erva come como o boi.
\VS{16}Eis que sua força está em seus lombos, e seu poder na musculatura de seu ventre.
\VS{17}Ele torna sua cauda dura como o cedro, e os nervos de suas coxas são entretecidos.
\VS{18}Seus ossos são
{\ADD{como}} tubos de bronze; seus membros, como barras de ferro.
\VS{19}Ele é a obra-prima dos caminhos de Deus; aquele que o fez o proveu de sua espada.
\FTNT{†}{{\FR 40:19 }
o proveu de sua espada
obscuro – trad. alt. é o que pode aproximar sua espada contra ele
}
\VS{20}Pois os montes lhe produzem pasto; por isso todos os animais do campo ali se alegram.
\FTNT{‡}{{\FR 40:20 }
alegram
ou brincam
}
\VS{21}Ele se deita debaixo das árvores sombrias; no esconderijo das canas e da lama.
\VS{22}As árvores sombrias o cobrem, cada uma com sua sombra; os salgueiros do ribeiro o cercam.
\VS{23}Ainda que o rio se torne violento, ele não se apressa; confia ainda que o Jordão transborde até sua boca.
\VS{24}Poderiam, por acaso, capturá-lo à vista de seus olhos,
{\ADD{ou}} com laços furar suas narinas?

\par }\Chap{41}{\PP \VerseOne{1}Poderás tu pescar ao leviatã com anzol, ou abaixar sua língua com uma corda?
\VS{2}Podes pôr um anzol em seu nariz, ou com um espinho furar sua queixada?
\VS{3}Fará ele súplicas a ti,
{\ADD{ou}} falará contigo suavemente?
\VS{4}Fará ele pacto contigo, para que tu o tomes por escravo perpétuo?
\VS{5}Brincarás tu com ele como com um passarinho, ou o atarás para tuas meninas?
\VS{6}Os companheiros farão banquete dele? Repartirão dele entre os mercadores?
\VS{7}Poderás tu encher sua pele de espetos, ou sua cabeça com arpões de pescadores?
\VS{8}Põe tua mão sobre ele; te lembrarás da batalha, e nunca mais voltarás a fazer.
\VS{9}Eis que a esperança de alguém
{\ADD{de vencê-lo}} falhará; pois, apenas ao vê-lo será derrubado.
\VS{10}Ninguém há
{\ADD{tão}} ousado que o desperte; quem pois,
{\ADD{ousa}} se opor a mim?
\VS{11}Quem me deu primeiro, para que eu
{\ADD{o}} recompense? Tudo o que há debaixo dos céus é meu.
\VS{12}Eu não me calarei a respeito de seus membros, nem de
{\ADD{suas}} forças, e da graça de sua estatura.
\VS{13}Quem descobrirá sua vestimenta superficial? Quem poderá penetrar sua couraça dupla?
\VS{14}Quem poderia abrir as portas de seu rosto? Ao redor de seus dentes há espanto.
\VS{15}Seus fortes escudos
\FTNT{*}{{\FR 41:15 }
escudos
i.e., provavelmente escamas
} são excelentes; cada um fechado, como um selo apertado.
\VS{16}Um está tão próximo do outro, que vento não pode entrar entre eles.
\VS{17}Estão grudados uns aos outros; estão tão travados entre si, que não se podem separar.
\VS{18}Cada um de seus roncos faz resplandecer a luz, e seus olhos são como os cílios do amanhecer.
\VS{19}De sua boca saem tochas, faíscas de fogo saltam dela.
\VS{20}De suas narinas sai fumaça, como de uma panela fervente ou de um caldeirão.
\VS{21}Seu fôlego acende carvões, e de sua boca sai chama.
\VS{22}A força habita em seu pescoço; diante dele salta-se de medo.
\VS{23}As dobras de sua carne estão apegadas
{\ADD{entre si}} ; cada uma está firme nele, e não podem ser movidas.
\VS{24}Seu coração é rígido como uma pedra, rígido como a pedra de baixo de um moinho.
\VS{25}Quando ele se levanta, os fortes tremem; por
{\ADD{seus}} abalos se recuam.
\VS{26}Se alguém lhe tocar com a espada, não poderá prevalecer; nem arremessar dardo, ou lança.
\VS{27}Ele considera o ferro como palha, e o aço como madeira podre.
\VS{28}A flecha não o faz fugir; as pedras de funda são para ele como sobras de cascas.
\VS{29}Considera toda arma como sobras de cascas, e zomba do mover da lança.
\VS{30}Por debaixo de si tem conchas pontiagudas; ele esmaga com suas pontas na lama.
\VS{31}Ele faz ferver as profundezas como a uma panela, e faz do mar como um pote de unguento.
\VS{32}Ele faz brilhar o caminho atrás de si; faz parecer ao abismo com cabelos grisalhos.
\VS{33}Não há sobre a terra algo que se possa comparar a ele. Ele foi feito para não temer.
\VS{34}Ele vê tudo que é alto; ele é rei sobre todos os filhos dos animais soberbos.
\FTNT{†}{{\FR 41:34 }
filhos dos animais soberbos
lit. filhos da soberba
}

\par }\Chap{42}{\PP \VerseOne{1}E Jó respondeu ao SENHOR, dizendo:
\VS{2}Eu sei que tudo podes, e nenhum de teus pensamentos pode ser impedido.
\VS{3}{\ADD{Tu dizes}} : Quem é esse que obscurece o conselho sem conhecimento? Por isso eu falei do que não entendia; coisas que eram maravilhosas demais para mim, e eu não as conhecia.
\VS{4}Escuta-me, por favor, e eu falarei; eu te perguntarei, e tu me ensina.
\VS{5}Com os ouvidos eu ouvira falar de ti; mas agora meus olhos te veem.
\VS{6}Por isso
{\ADD{me}} abomino, e me arrependo no pó e na cinza.
\VS{7}E sucedeu que, depois que o SENHOR acabou de falar essas palavras a Jó, o SENHOR disse a Elifaz, o temanita: Minha ira se acendeu contra ti e contra os teus dois amigos; porque não falastes de mim o que era correto, como o meu servo Jó.
\VS{8}Por isso tomai sete bezerros e sete carneiros, ide ao meu servo Jó, e oferecei holocaustos em vosso favor, então meu servo Jó orará por vós; pois certamente atenderei a ele para eu não vos tratar conforme a vossa insensatez; pois não falastes de mim o que era correto, como o meu servo Jó.
\VS{9}Então foram Elifaz, o temanita, Bildade, o suíta, e Zofar, o naamita, e fizeram como o SENHOR havia lhes dito; e o SENHOR atendeu a Jó.
\VS{10}E o SENHOR livrou Jó de seu infortúnio, enquanto ele orava pelos seus amigos; e o SENHOR acrescentou a Jó o dobro de tudo quanto ele{\ADD{ antes}} possuía.
\VS{11}Então vieram e ele todos os seus irmãos, e todas as suas irmãs, e todos os que o conheciam antes; e comeram com ele pão em sua casa, e condoeram-se dele, e o consolaram acerca de toda a calamidade
\FTNT{*}{{\FR 42:11 }
calamidade
lit. mal
} que o SENHOR tinha trazido sobre ele; e cada um deles lhe deu uma peça de dinheiro, e uma argola de ouro.
\VS{12}E
{\ADD{assim}} o SENHOR abençoou o último estado de Jó mais que o seu primeiro; porque teve catorze mil ovelhas, seis mil camelos, mil juntas de bois, e mil asnas.
\VS{13}Também teve sete filhos e três filhas.
\VS{14}E chamou o nome da uma Jemima, e o nome da segunda Quézia, e o nome da terceira Quéren-Hapuque.
\VS{15}E em toda a terra não se acharam mulheres tão belas como as filhas de Jó; e seu pai lhes deu herança entre seus irmãos.
\VS{16}E depois disto Jó viveu cento e quarenta anos; e viu seus filhos, e os filhos de seus filhos, até quatro gerações.
\VS{17}Então Jó morreu, velho, e farto de dias.\par }